\documentclass[11pt,a4paper]{article}
\usepackage[T1]{fontenc}
\usepackage[utf8]{inputenc}
\usepackage{lmodern,microtype}
\usepackage[margin=25mm,headheight=14pt]{geometry}
\usepackage{amsmath,amssymb,amsthm,mathtools}
\usepackage{booktabs,array,enumitem,xcolor,fancyhdr,listings,textcomp}
\usepackage[colorlinks=true,linkcolor=blue!45!black,citecolor=blue!45!black,urlcolor=blue!45!black]{hyperref}
\hypersetup{pdftitle={The Randic index is at least the average distance: a proof of Fajtlowicz's conjecture},
 pdfsubject={Randic index, average distance, Graffiti conjecture, weighted distance inequality, stationarity-deletion principle}}
\pagestyle{fancy}\fancyhf{}
\fancyhead[L]{\small The Randi\'c index is at least the average distance}
\fancyhead[R]{\small VMRC research manuscript}
\fancyfoot[C]{\thepage}
\newtheorem{theorem}{Theorem}[section]
\newtheorem{lemma}[theorem]{Lemma}
\newtheorem{proposition}[theorem]{Proposition}
\newtheorem{corollary}[theorem]{Corollary}
\theoremstyle{definition}
\newtheorem{definition}[theorem]{Definition}
\theoremstyle{remark}
\newtheorem{remark}[theorem]{Remark}
\numberwithin{equation}{section}
\newcommand{\ad}{\operatorname{ad}}
\newcommand{\diam}{\operatorname{diam}}
\newcommand{\supp}{\operatorname{supp}}
\newcommand{\hi}{\operatorname{hi}}
\lstset{basicstyle=\ttfamily\scriptsize,breaklines=true,columns=fullflexible,keepspaces=true,
  frame=single,framerule=0.3pt,rulecolor=\color{black!40},showstringspaces=false,
  upquote=true}
\setlength{\emergencystretch}{2em}

\title{\bfseries The Randi\'c index is at least the average distance:\\ a proof of Fajtlowicz's conjecture\\[5pt]
\large\normalfont A stationarity--deletion principle for weighted distance sums}
\author{Verifiable Mathematical Reasoning Corpus --- research manuscript\thanks{Prepared with an AI system (Claude, Anthropic) as a continuation of the Verifiable Mathematical Reasoning Corpus (VMRC). The programs for the finite verification are reproduced in Appendix~\ref{app:code}. This manuscript has not been refereed.}}
\date{20 November 2026}

\begin{document}
\maketitle

\begin{center}
\fbox{\parbox{0.93\linewidth}{\small\textbf{Status.}
This manuscript claims a complete proof of Fajtlowicz's 1988 conjecture
$R(G)\ge\ad(G)$. One step of the proof is a finite exhaustive computation over
the pendant-vertex decorations of all two-connected graphs with at most nine
vertices (Section~\ref{sec:budget} and Appendix~\ref{app:code}): the combinatorial
quantities are computed in exact integer arithmetic, and the final comparison is
made in floating point with a margin of $0.38$. It was carried out by two
independently written programs. Everything else is proved in the text. The argument
was produced with an AI system and checked by the computations listed in
Section~\ref{sec:remarks} and by a separate AI review pass, but it has not been
refereed by humans or formally verified, and priority cannot be certified by a
literature search alone (Section~\ref{sec:audit}).}}
\end{center}

\begin{abstract}
For a connected graph $G$ let $R(G)=\sum_{uv\in E(G)}(d_ud_v)^{-1/2}$ be its
Randi\'c index and $\ad(G)$ the average distance between two distinct vertices.
In 1988 Fajtlowicz conjectured, on the basis of the program Graffiti, that
$R(G)\ge\ad(G)$ for every connected graph. The conjecture was known for trees
(Cygan, Pilipczuk and \v{S}krekovski), for unicyclic graphs, and under various
restrictions on degrees, diameter or blocks. We prove it for all connected graphs,
with equality only for $K_2$.

The proof goes through the \emph{weighted distance inequality}
\[
 \sum_{x\ne y}p_xp_y\,d_G(x,y)+2\sum_xp_x^2\ \le\ M(G)+1,\qquad
 M(G)=\sum_{uv\in E(G)}\frac1{\max(d_u,d_v)},
\]
valid for every probability vector $p$ on the vertices of a connected graph with
at least two vertices. Two-point measures give the diameter theorem
$M(G)\ge\diam(G)/2$ of Dvo\v{r}\'ak, Lidick\'y and \v{S}krekovski, and the uniform
measure, combined with an elementary comparison of $R$ and $M$, gives Fajtlowicz's
conjecture. The new tool is a
\emph{stationarity--deletion principle}. At a maximiser $p$ of the left-hand side
the first-order optimality conditions turn the deletion of a vertex $u$ of mass
$t=p_u<\frac12$ into the inequality
$(1-t)^2\bigl(M(G)-M(G-u)\bigr)<(M(G)-1)\,t^2$ for a minimal counterexample. Thus
a vertex whose deletion lowers $M$ has to carry a definite amount of mass, and the
total mass one becomes a budget. For the graphs that survive the classical reductions
(two-connected cores with pendant vertices, all of whose core edges are
``supportive'') we show that the budget is exceeded: by an analytic estimate when
the core has at least ten vertices, and by an exhaustive computer enumeration of
the decorations of the $201{,}727$ two-connected graphs with at most nine vertices,
carried out by two independently written programs.
\end{abstract}

\tableofcontents

\section{Introduction}\label{sec:intro}

\subsection{The conjecture}
All graphs are finite, simple and undirected. For a graph $G$ write $d_v=d_G(v)$
for the degree of $v$, and for a connected graph of order $n\ge2$ let $d_G(x,y)$
be the distance and
\[
 R(G)=\sum_{uv\in E(G)}\frac1{\sqrt{d_ud_v}},\qquad
 W(G)=\sum_{\{x,y\}}d_G(x,y),\qquad \ad(G)=\frac{W(G)}{\binom n2}.
\]
The index $R$ was introduced by Randi\'c \cite{Randic}; $\ad$ is the average
(mean) distance. In his report on the conjectures of the program Graffiti,
Fajtlowicz \cite{Faj} proposed two inequalities linking $R$ with distances:
$R(G)\ge\operatorname{rad}(G)-1$, and the following, which is the subject of
this paper (we quote it in the form used in \cite{CPS,YL,LiShi}).

\medskip\noindent\textbf{Conjecture (Fajtlowicz, 1988).}
\emph{For every connected graph $G$, $R(G)\ge\ad(G)$.}
\medskip

Cygan, Pilipczuk and \v{S}krekovski \cite{CPS} proved the conjecture for trees, You
and Liu \cite{YL} for unicyclic graphs, Li and Shi \cite{LiShi} for graphs with
$\delta(G)\ge n/5$ and $n\ge15$, and Dalf\'o \cite{Dalfo} for graphs of sufficiently
large order whose minimum degree exceeds (approximately) $\Delta^{1/3}$. The
unrefereed collection \cite{VMRC189} of the Verifiable Mathematical Reasoning
Corpus contains nine manuscripts on the conjecture (its papers 06--14). They prove
it for diameter at most $53$, for maximum degree at most $100$, for circuit rank at
most $155$, and for graphs all of whose blocks satisfy certain diameter or degree
conditions. Papers 11--14 of that collection introduce the weighted inequality of
Theorem~\ref{thm:W} below, prove it when all distances are truncated at $40$, and
reduce it to two-connected graphs with pendant vertices; the collection records
that the unrestricted conjecture remained open there. The first of the two
conjectures is related to the theorem $M(G)\ge\diam(G)/2$ of Dvo\v{r}\'ak,
Lidick\'y and \v{S}krekovski \cite{DLS}, where
\[
 M(G)=\sum_{uv\in E(G)}\frac1{\max(d_u,d_v)}\ \le\ R(G)
\]
is their modified index (denoted $R'$ in \cite{DLS}).

\subsection{Results}
For a connected graph $G$ on vertex set $V$ and a probability vector $p$ on $V$
put
\begin{equation}\label{eq:F}
 F(G;p)=\sum_{x\ne y}p_xp_y\,d_G(x,y)+2\sum_{x}p_x^2,\qquad F(G)=\max_pF(G;p),
\end{equation}
the first sum being over ordered pairs of distinct vertices. The maximum exists
because $F(G;\cdot)$ is continuous on the simplex of probability vectors; a vector
at which it is attained is called a \emph{maximiser}.

\begin{theorem}[Weighted distance inequality]\label{thm:W}
For every connected graph $G$ with at least two vertices, $F(G)\le M(G)+1$.
\end{theorem}

\begin{theorem}[Fajtlowicz's conjecture]\label{thm:F}
For every connected graph $G$ with at least two vertices, $R(G)\ge\ad(G)$, with
equality if and only if $G=K_2$.
\end{theorem}

\begin{corollary}\label{cor:main}
Let $G$ be a connected graph of order $n\ge2$.
\begin{enumerate}[label=(\alph*)]
\item {\rm (Dvo\v{r}\'ak--Lidick\'y--\v{S}krekovski \cite{DLS})} $M(G)\ge\diam(G)/2$.
\item For any $k\ge2$ distinct vertices $x_1,\dots,x_k$,
$\ \frac1{k^2}\sum_{i\ne j}d_G(x_i,x_j)\le M(G)+1-\frac2k$.
\item $M(G)\ge\bigl(1-\frac1n\bigr)\ad(G)-1+\frac2n$.
\end{enumerate}
\end{corollary}

Theorem~\ref{thm:W} is sharp in several ways: equality holds for every star with
at least two leaves and every measure carried by its leaves, for every path with at
least three vertices and the two-point measure on its ends, and for some graphs
with cycles (Remark~\ref{rem:sharp}).

\subsection{Method}
Theorem~\ref{thm:F} follows from Theorem~\ref{thm:W} by an elementary comparison of
$R$ and $M$ (Section~\ref{sec:proofs}). Theorem~\ref{thm:W} is proved by considering a
counterexample $G$ with the fewest vertices and then the fewest edges, together
with a maximiser $p$ of $F(G;\cdot)$.

\emph{Classical reductions} (Section~\ref{sec:red}). Deleting a cycle edge that does
not increase $M$, and the cut-vertex padding of \cite[Lemma~4]{DLS} in the weighted
form of \cite{VMRC189}, reduce $G$ to a two-connected core with at most one pendant
vertex at each core vertex, in which the deletion of any core edge strictly
increases $M$ (an \emph{irreducible decorated block}).

\emph{Stationarity--deletion} (Section~\ref{sec:sd}; this is the new ingredient). Let
$u$ be a vertex whose deletion leaves a connected graph, with mass $t=p_u$.
Because $p$ is a maximiser, $\sum_yp_y\,d(u,y)=F(G)-2t$ if $t>0$. Deleting $u$,
renormalising $p$ and applying minimality to $G-u$ therefore gives a closed
inequality in which the unknown distances have disappeared:
\[
 (1-2t)\bigl(F(G)-M(G)-1\bigr)\ \le\ (M(G)-1)\,t^2-(1-t)^2\bigl(M(G)-M(G-u)\bigr).
\]
In a counterexample the left side is positive when $t<\frac12$, so
$t\ge\theta\bigl(M(G)-M(G-u)\bigr)$ for an explicit increasing function $\theta$:
\emph{a vertex whose deletion lowers $M$ must carry mass}. Summing over suitable
deletable vertices shows that a minimal counterexample satisfies $\Theta(G)\le1$
for a purely combinatorial quantity $\Theta$ (the \emph{mass budget}).

\emph{The budget is exceeded} (Section~\ref{sec:budget}). We prove $\Theta(G)>1$ for
every irreducible decorated block. If the core has $h\ge10$ vertices this follows
from a lower bound for the sum of the deletion values of the core vertices, the
constraints that supportive edges impose on them, and a counting bound
$M\ge\frac{15}4$; for $h\le9$ it is verified by an exhaustive computer enumeration.

\paragraph{New problem, new method, application.}
The paper has the following logical structure. \emph{(i) A new problem}: the
\emph{budget problem}, which asks whether $\Theta(G)>1$ for every irreducible
decorated block $G$ (Proposition~\ref{prop:exceed}); it is a statement about degrees
only and does not mention distances or measures. \emph{(ii) A new method}: the
stationarity--deletion principle (Lemma~\ref{lem:sdel} and
Proposition~\ref{prop:budget}), which converts any minimal counterexample to the
weighted inequality into a negative instance of the budget problem, and the
site-deletion calculus of Section~\ref{sec:blocks}, which solves the budget
problem. \emph{(iii) Application to a known problem}: Theorem~\ref{thm:W} and,
through it, Fajtlowicz's conjecture in its original form (Theorem~\ref{thm:F}).

\paragraph{Relation to earlier work.}
The functional \eqref{eq:F} and the inequality of Theorem~\ref{thm:W} were
introduced, as an open problem, in papers 11--14 of the collection \cite{VMRC189}.
Those papers also contain the deduction of Fajtlowicz's inequality from the weighted
one (reproduced in Section~\ref{sec:proofs}), the weighted form of the cut-vertex
padding of \cite{DLS} (Lemma~\ref{lem:cut}), and the description of an edge-minimal
obstruction as a two-connected core with at most one pendant vertex at each core
vertex in which the deletion of any core edge strictly increases $M$, together with
the consequences (P1) and (P2) of Lemma~\ref{lem:struct}. Lemma~\ref{lem:cut} has the
same statement as Lemma~4 of \cite{DLS}; the use of vertices without neighbours of
smaller degree goes back to Lemma~5 of \cite{DLS}, and the counting argument of
Lemma~\ref{lem:count} to Lemma~7 of \cite{DLS} (numbering of the preprint version).
We give complete proofs of all these statements. Thus no result of the unrefereed
collection \cite{VMRC189} is used as a black box, and neither is the theorem of
\cite{DLS}, which is reproved as Corollary~\ref{cor:main}(a). The only published
result used is the bound $W(G)\le\binom{n+1}3$ \cite{EJS}, for which a proof is
included.

New in this paper are the stationarity--deletion inequality
(Lemma~\ref{lem:sdel}), the mass budget (Proposition~\ref{prop:budget}), the
deletion values of the sites with their sum inequality and upper bounds
(Lemmas~\ref{lem:X}--\ref{lem:upper}), and the proof that the budget is exceeded
(Proposition~\ref{prop:exceed}), including the enumeration behind
Lemma~\ref{lem:small}. In \cite{VMRC189} the weighted inequality had been verified
directly for all cores with at most eight vertices and proved for truncated
distances; the first-order optimality conditions of the maximiser were not used
there.

\section{Notation and elementary formulas}\label{sec:prelim}

Let $G$ be a graph. For an edge $uv$ its \emph{weight} is $1/\max(d_u,d_v)$, so
$M(G)$ is the total weight. If $G$ has an edge, then $M(G)\ge1$: the $\Delta$ edges
at a vertex of maximum degree $\Delta$ have weight $1/\Delta$ each. For a vertex $w$
with $d_w\ge2$ put
\[
 g_w=d_w(d_w-1),\qquad \ell(w)=\#\{z\in N(w):d_z<d_w\},
\]
the number of \emph{lower} neighbours of $w$. A \emph{leaf} is a vertex of degree
one. For a statement $S$ we write $[S]$ for $1$ if $S$ holds and $0$ otherwise.
For a connected graph $G$ let $K_G$ be the symmetric matrix with $K_G(x,y)=d_G(x,y)$
for $x\ne y$ and $K_G(x,x)=2$, so that $F(G;p)=p^{\mathsf T}K_Gp$.

\begin{lemma}[Monotonicity]\label{lem:mono}
Let $G'$ be a connected subgraph of a connected graph $G$ and let $p$ be a
probability vector on $V(G)$ with $\supp p\subseteq V(G')$. Then
$F(G;p)\le F(G';p)$.
\end{lemma}
\begin{proof}
$d_{G'}(x,y)\ge d_G(x,y)$ for $x,y\in V(G')$, and the diagonal entries agree.
\end{proof}

\begin{lemma}[Deletion formulas]\label{lem:del}
\begin{enumerate}[label=(\alph*)]
\item Let $uv\in E(G)$ with $d_u,d_v\ge2$. Then
\[
 M(G-uv)-M(G)=\frac{\ell(u)-[d_v<d_u]}{g_u}+\frac{\ell(v)-[d_u<d_v]}{g_v}
 -\frac1{\max(d_u,d_v)} .
\]
\item Let $v'$ be a leaf with neighbour $v$, where $d_v\ge2$. Then
$M(G)-M(G-v')=\dfrac{d_v-\ell(v)}{g_v}\ \ge0$.
\item Let $v\in V(G)$, let $D_0$ be a set of leaves adjacent to $v$, put
$D=\{v\}\cup D_0$ and $N=N(v)\setminus D_0$, and assume $d_w\ge2$ for all $w\in N$.
Then
\[
 M(G)-M(G-D)=\sum_{w\in N(v)}\frac1{\max(d_v,d_w)}
 -\sum_{w\in N}\frac{\ell(w)-[d_v<d_w]}{g_w}-\tau,
\]
where $\tau\ge0$ is the sum of $1/g_w$ over all edges $wz$ of $G$ with $w,z\in N$
and $d_w=d_z$.
\end{enumerate}
\end{lemma}
\begin{proof}
If the degree of $w$ drops from $d_w$ to $d_w-1$ while the degree of a neighbour
$z$ is unchanged, the weight of $wz$ changes exactly when $d_z<d_w$, and then it
increases from $1/d_w$ to $1/(d_w-1)$, that is, by $1/g_w$. If the degrees of both
ends of an edge drop by one, its weight increases by $1/g$ of the end of larger
degree.

(a) Only $d_u$ and $d_v$ change. The edges at $u$ other than $uv$ whose weight
increases are those to the $\ell(u)-[d_v<d_u]$ lower neighbours of $u$ different
from $v$; similarly at $v$; and the weight $1/\max(d_u,d_v)$ of $uv$ is lost.

(b) The lost edge has weight $1/d_v$, and the weight increases on the $\ell(v)-1$
edges from $v$ to its lower neighbours other than $v'$. Hence the difference is
$1/d_v-(\ell(v)-1)/g_v=(d_v-\ell(v))/g_v$, which is nonnegative because
$\ell(v)\le d_v$.

(c) The lost edges are those at $v$ (a leaf edge $vv'$ has weight $1/d_v$). Exactly
the vertices of $N$ lose one unit of degree. For $w\in N$ and a neighbour
$z\notin N\cup D$ of $w$ the weight of $wz$ increases by $1/g_w$ iff $d_z<d_w$.
For an edge $wz$ with $w,z\in N$ the weight increases by $1/g_w$ if $d_z<d_w$,
by $1/g_z$ if $d_w<d_z$, and by $1/g_w$ if $d_w=d_z$. Summing, every $w\in N$
contributes $1/g_w$ for each of its lower neighbours other than $v$, and the edges
inside $N$ with equal end degrees contribute $\tau$.
\end{proof}

\section{Reductions}\label{sec:red}

Call a connected graph $G$ with at least two vertices \emph{good} if
$F(G)\le M(G)+1$ and \emph{bad} otherwise. Theorem~\ref{thm:W} says that there are
no bad graphs. \emph{Throughout Sections~\ref{sec:red} and \ref{sec:sd} we assume
that bad graphs exist, we let $G$ be a bad graph for which the pair
$(|V(G)|,|E(G)|)$ is lexicographically smallest, and we fix a maximiser $p$ of
$F(G;\cdot)$.} Thus $F(G;p)=F(G)>M(G)+1$, and every connected graph with at least
two vertices that has fewer vertices than $G$, or the same number of vertices and
fewer edges, is good.

\begin{lemma}\label{lem:base}
$K_2$, all stars and all double stars are good. (A double star consists of two
adjacent vertices, each with at least one further neighbour, all further vertices
being leaves.)
\end{lemma}
\begin{proof}
Let $c$ be the centre of a star $K_{1,k}$, $k\ge1$. Then $K(x,y)=2$ unless
$\{x,y\}$ consists of $c$ and a leaf, where it is $1$; hence
$F(G;p)=2-2p_c(1-p_c)\le2=M(G)+1$. In a double star with centres $a,b$ write
$\alpha,\beta$ for the masses of the leaves at $a$ and at $b$. Since $K\le2$
except for $K=3$ between leaves at different centres, and $K=1$ on the pairs
$\{a,\text{leaf of }a\}$, $\{b,\text{leaf of }b\}$, $\{a,b\}$,
\[
 F(G;p)=2+2\alpha\beta-2p_a\alpha-2p_b\beta-2p_ap_b\le2+2\alpha\beta\le\tfrac52 .
\]
If $a,b$ have $k\ge l\ge1$ leaves, then $M(G)=\frac{k}{k+1}+\frac{l}{l+1}+\frac1{k+1}
=1+\frac l{l+1}\ge\frac32$.
\end{proof}

\begin{lemma}[Subgraph principle]\label{lem:sub}
Let $G'$ be a connected proper subgraph of $G$ with at least two vertices such that
$\supp q\subseteq V(G')$ for some maximiser $q$ of $F(G;\cdot)$. Then
$M(G')>M(G)$. In particular:
\begin{enumerate}[label=(\alph*)]
\item every edge $e$ of $G$ that lies on a cycle is \emph{supportive}:
$M(G-e)>M(G)$;
\item every leaf of $G$ has positive mass under every maximiser of $F(G;\cdot)$.
\end{enumerate}
\end{lemma}
\begin{proof}
$G'$ has fewer vertices than $G$, or the same vertices and fewer edges, so it is
good. By Lemma~\ref{lem:mono}, $M(G)+1<F(G;q)\le F(G';q)\le M(G')+1$. For (a),
$G-e$ is connected and contains all vertices. For (b), suppose $q_z=0$ for a leaf $z$ and a
maximiser $q$. As $G\ne K_2$ by Lemma~\ref{lem:base}, the neighbour of $z$ has degree
at least $2$ and $G'=G-z$ is connected with at least two vertices; then
$M(G')>M(G)$ contradicts Lemma~\ref{lem:del}(b).
\end{proof}

\begin{lemma}[Cut vertices]\label{lem:cut}
For every cut vertex $v$ of $G$, all components of $G-v$ except one consist of a
single vertex.
\end{lemma}
\begin{proof}
Suppose two components of $G-v$ have more than one vertex. Then
$V(G)\setminus\{v\}=A\sqcup B$ with no edges between $A$ and $B$, where each of $A$
and $B$ contains such a component. Let $G_A$ be obtained from $G[A\cup\{v\}]$ by
attaching $d_B:=|N(v)\cap B|$ new leaves to $v$, and define $G_B$ symmetrically with
$d_A=|N(v)\cap A|$ new leaves. Here $d_A,d_B\ge1$ because $G$ is connected, and
$G_A$, $G_B$ are connected graphs with at least two vertices. All vertices of $A\cup\{v\}$ have the same degree in
$G_A$ as in $G$, so the old edges keep their weights and the new edges have total
weight $d_B/d_v$; hence
\begin{equation}\label{eq:padM}
 M(G_A)+M(G_B)=M(G)+\frac{d_A+d_B}{d_v}=M(G)+1 .
\end{equation}
$G_A$ has $|A|+1+d_B\le|V(G)|$ vertices. If equality holds, every vertex of $B$ is
adjacent to $v$, and $G_A$ has $|E(G[B])|$ fewer edges than $G$; this number is
positive because $B$ contains a component with an edge. So $G_A$, and likewise
$G_B$, is good.

Fix new leaves $a^\ast$ of $G_A$ and $b^\ast$ of $G_B$ and define
$\pi_A\colon V(G)\to V(G_A)$ by $\pi_A(x)=x$ for $x\in A$ and $\pi_A(x)=a^\ast$
otherwise; define $\pi_B$ symmetrically. We claim that for all $x,y\in V(G)$
\begin{equation}\label{eq:padK}
 K_{G_A}(\pi_Ax,\pi_Ay)+K_{G_B}(\pi_Bx,\pi_By)\ \ge\ K_G(x,y)+2 .
\end{equation}
Indeed, distances inside $A\cup\{v\}$ are the same in $G$ and $G_A$. If
$x,y\in A$, the left side is $K_G(x,y)+K_{G_B}(b^\ast,b^\ast)=K_G(x,y)+2$. If
$x\in A$ and $y\in B$, it is $(d(x,v)+1)+(d(v,y)+1)=d_G(x,y)+2$. If $x\in A$ and
$y=v$, it is $(d(x,v)+1)+2=K_G(x,v)+3$. If $x=y=v$, it is $2+2=K_G(v,v)+2$. The
remaining cases follow by exchanging $x$ and $y$, or $A$ and $B$. Multiplying \eqref{eq:padK} by $p_xp_y$ and summing gives
$F(G_A;\pi_Ap)+F(G_B;\pi_Bp)\ge F(G;p)+2$ for the image measures. As $G_A,G_B$ are
good, \eqref{eq:padM} yields $F(G;p)\le M(G_A)+1+M(G_B)+1-2=M(G)+1$, a contradiction.
\end{proof}

\begin{definition}\label{def:block}
A \emph{decorated block} is a connected graph $G$ such that, $L$ being the set of
its leaves, $H=G-L$ is two-connected with $h=|V(H)|\ge3$ vertices and every vertex
of $H$ is adjacent to at most one leaf. The vertices of $H$ are the \emph{sites},
its edges the \emph{core edges}; a site $v$ is \emph{loaded} if it has a leaf,
denoted $v'$, and \emph{unloaded} otherwise. A decorated block is
\emph{irreducible} if every core edge $e$ is supportive, i.e.\ $M(G-e)>M(G)$.
\end{definition}

\begin{proposition}\label{prop:shape}
The minimal bad graph $G$ is an irreducible decorated block, and $p_v=0$ for every
loaded site $v$.
\end{proposition}
\begin{proof}
Let $L$ be the set of leaves and $H=G-L$. Since $G\ne K_2$ (Lemma~\ref{lem:base}),
$H$ is a non-empty connected graph. If $|V(H)|=1$ then $G$ is a star, and if $|V(H)|=2$ then $G$ is a double star (both vertices of $H$ are
non-leaves, so each has a leaf); by Lemma~\ref{lem:base} neither occurs, so
$h=|V(H)|\ge3$. Suppose $v$ is a cut vertex of $H$ and let $C_1,C_2$ be two
components of $H-v$. The component of $G-v$ containing $C_i$ contains a vertex
$w\in V(H)$; it is not the single vertex $w$, because then $w$ would be a leaf of
$G$. This contradicts Lemma~\ref{lem:cut}. Hence $H$ is two-connected.

Suppose a site $s$ has two leaves $a,b$. The rows of $K_G$ at $a$ and $b$ agree
outside $\{a,b\}$ and $K(a,a)=K(b,b)=K(a,b)=2$; so moving the mass of $b$ to $a$
does not change $F$ and produces a maximiser $q$ with $q_b=0$, contradicting
Lemma~\ref{lem:sub}(b). Thus $G$ is a decorated block. Core edges lie on cycles of
$H$, so $G$ is irreducible by Lemma~\ref{lem:sub}(a).

Finally let $v$ be a loaded site and $t=p_v$. Since $K(v',y)=K(v,y)+1$ for
$y\notin\{v,v'\}$, $K(v',v)=K(v,v)-1$ and $K(v',v')=K(v,v')+1$, we get
$(Kp)_{v'}-(Kp)_v=1-2t$, and therefore
$F(G;p+t(\delta_{v'}-\delta_v))-F(G;p)=2t(1-2t)+t^2(2+2-2)=2t(1-t)$. As $p$ is a
maximiser, $t\in\{0,1\}$, and $t=1$ is impossible because then $F(G;p)=2\le M(G)+1$.
\end{proof}

\section{Irreducible decorated blocks}\label{sec:blocks}

In this section $G$ is an arbitrary irreducible decorated block with core $H$ on
$h$ sites. Degrees $d_v$ are always taken in $G$. For a site $v$ let $\epsilon_v\in\{0,1\}$
indicate whether it is loaded, and let $\ell(v)$, $e(v)$, $\hi(v)$ be the numbers of
neighbours of $v$ in $G$ of smaller, equal and larger degree (the leaf $v'$ is a
lower neighbour). Put
\[
 a(v)=\frac{\ell(v)}{g_v},\qquad b(v)=\frac{d_v-\ell(v)}{g_v}=\frac{e(v)+\hi(v)}{g_v},
 \qquad a(v)+b(v)=\frac1{d_v-1}.
\]

\begin{lemma}[Structure]\label{lem:struct}
\begin{enumerate}[label=\textup{(P\arabic*)}]
\item Every site has a neighbour of strictly smaller degree: $\ell(v)\ge1$.
\item Every site has degree at least $3$. A site of degree $3$ is loaded and its
two core neighbours have degree at least $4$.
\item For every core edge $uv$: if $d_u<d_v$ then $a(u)>b(v)$; if $d_u=d_v$ then
$a(u)\ge b(v)$ and $a(v)\ge b(u)$.
\end{enumerate}
\end{lemma}
\begin{proof}
Sites have at least two core neighbours, so Lemma~\ref{lem:del}(a) applies to core
edges. Let $uv$ be a core edge with $d_u\le d_v$. If $d_u<d_v$, supportiveness reads
$\ell(u)/g_u+(\ell(v)-1)/g_v>1/d_v=(d_v-1)/g_v$, i.e.\ $a(u)>b(v)$. If $d_u=d_v=d$
it reads $\ell(u)+\ell(v)>d-1$, i.e.\ $\ell(u)\ge d-\ell(v)$, which is $a(u)\ge b(v)$,
and symmetrically. This is (P3). If $\ell(u)=0$ for a site $u$, take a core
neighbour $v$; then $d_v\ge d_u$ and $a(u)=0\ge b(v)\ge0$ forces $d_u=d_v$ and
$b(v)=0$, i.e.\ $\ell(v)=d_v$, which is impossible because $u$ is not a lower
neighbour of $v$. This proves (P1). A site of minimum degree has a lower neighbour,
which must be its leaf; since it also has two core neighbours, all sites have degree
at least $3$. Let $v$ be a site of degree $3$. Its lower neighbour has degree at most
$2$, hence is its leaf, so $v$ has exactly two core neighbours and $\ell(v)=1$. If a
core neighbour $w$ had degree $3$ as well, then $\ell(v)+\ell(w)=2<3$ would
contradict (P3).
\end{proof}

A site $v$ is a \emph{strict peak} if all its core neighbours have smaller degree.
Strict peaks are pairwise non-adjacent, and a loaded strict peak has degree at
least $4$.

\paragraph{Deletion values.}
For a site $v$ define
\[
 X(v)=M(G)-M\bigl(G-\{v,v'\}\bigr)\quad(\text{$v'$ omitted if $v$ is unloaded}),
\]
and note that $G-\{v,v'\}$ is connected and has at least two vertices, because
$H-v$ is connected. For a loaded site, $G-v'$ is connected as well, and by
Lemma~\ref{lem:del}(b)
\begin{equation}\label{eq:bleaf}
 M(G)-M(G-v')=b(v)\ \ge0 .
\end{equation}

\begin{lemma}\label{lem:X}
For every site $v$,
\[
 X(v)=\frac{\epsilon_v}{d_v}+\sum_{w\in N_H(v)}c(v,w)-\tau_v,\qquad
 c(v,w)=\begin{cases} b(w),& d_w>d_v,\\[2pt] b(w)-\dfrac1{g_w},& d_w=d_v,\\[6pt]
 \dfrac1{d_v}-a(w),& d_w<d_v,\end{cases}
\]
where $\tau_v\ge0$ is the sum of $1/g_w$ over the core edges $wz$ with
$w,z\in N_H(v)$ and $d_w=d_z$. Moreover $X(v)\le1$.
\end{lemma}
\begin{proof}
Apply Lemma~\ref{lem:del}(c) with $D_0=\{v'\}$ or $\emptyset$. The term of
$w\in N_H(v)$ is $1/\max(d_v,d_w)-(\ell(w)-[d_v<d_w])/g_w$. If $d_w>d_v$ this is
$(d_w-1-\ell(w)+1)/g_w=b(w)$; if $d_w=d_v$ it is $(d_w-1-\ell(w))/g_w=b(w)-1/g_w$;
if $d_w<d_v$ it is $1/d_v-a(w)$. Finally, the terms subtracted in
Lemma~\ref{lem:del}(c) are non-negative (if $d_v<d_w$ then $v$ is a lower neighbour
of $w$, so $\ell(w)\ge1$). Hence $X(v)$ is at most the total weight of the edges at
$v$, which is at most $d_v\cdot\frac1{d_v}$.
\end{proof}

Let $M_{\rm core}$ be the total weight of the core edges, so that
$M(G)=M_{\rm core}+\sum_{v\ \rm loaded}1/d_v$.

\begin{lemma}[Sum inequality]\label{lem:sum}
$\displaystyle\sum_{v\in V(H)}X(v)\ \ge\ M_{\rm core}$.
\end{lemma}
\begin{proof}
$\sum_vX(v)$ equals the sum over the sites $v$ of the weight of the edges at $v$
(including $vv'$), which is $2M_{\rm core}+\sum_{\rm loaded}1/d_v$, minus the sum
over $v$ of the total increase of the weights of the other edges when
$\{v,v'\}$ is deleted. We compute the latter edge by edge, using the proof of
Lemma~\ref{lem:del}. A leaf edge $yy'$ increases by $1/g_y$ for each of the
$d_y-1$ core neighbours $v$ of $y$: in total $1/d_y$. A core edge $xy$ with
$d_x<d_y$ increases by $1/g_y$ exactly when $v$ is a core neighbour of $y$ other
than $x$: in total $(d_y-1-\epsilon_y)/g_y\le 1/d_y$, its weight. A core edge $xy$ with
$d_x=d_y=d$ increases by $1/g_x$ exactly when $v$ is a common core neighbour of
$x$ and $y$; there are at most $d-1$ of these, so the total is at most $1/d$, its
weight. Hence the total increase is at most
$\sum_{\rm loaded}1/d_y+M_{\rm core}$.
\end{proof}

\begin{lemma}[Upper bounds]\label{lem:upper}
Let $v$ be a loaded site of degree $d$ with $n_l$, $n_e$, $n_h$ core neighbours
of smaller, equal and larger degree $(n_l+n_e+n_h=d-1)$. Then
\begin{equation}\label{eq:U}
 X(v)\ \le\ \frac{(n_l+1)(d-1)+n_h}{d(d-1)} .
\end{equation}
Consequently, with $\xi(v):=\max\{X(v),b(v)\}$:
\begin{enumerate}[label=(\roman*)]
\item if $d=3$, then $\frac13\le\xi(v)\le\frac23$;
\item if $d\ge4$ and $v$ is not a strict peak, then
$\dfrac1{d(d-1)}\le\xi(v)\le1-\dfrac{d-2}{d(d-1)}$;
\item if $v$ is a strict peak of degree $4$, then $\xi(v)=\frac12$.
\end{enumerate}
\end{lemma}
\begin{proof}
By (P3), $c(v,w)=b(w)\le a(v)$ if $d_w>d_v$; $c(v,w)=b(w)-1/g_w\le a(v)-1/g_v$ if
$d_w=d_v$; and $c(v,w)=1/d_v-a(w)\le1/d_v-b(v)$ if $d_w<d_v$. Here
$a(v)=(n_l+1)/g_v$ and $b(v)=(n_e+n_h)/g_v$. With $g=g_v$, Lemma~\ref{lem:X} gives
\[
 gX(v)\le(d-1)+n_h(n_l+1)+n_e n_l+n_l\bigl((d-1)-n_e-n_h\bigr)
 =(n_l+1)(d-1)+n_h ,
\]
which is \eqref{eq:U}. Also $b(v)=(n_e+n_h)/g$ does not exceed the right side of
\eqref{eq:U}. (i) By (P2), $n_l=n_e=0$ and $n_h=2$, so $b(v)=\frac13$ and the bound
is $\frac46$. (ii) Here $n_e+n_h\ge1$, so $\xi(v)\ge b(v)\ge1/g$, and
$(n_l+1)(d-1)+n_h\le n_l(d-2)+2(d-1)\le(d-2)^2+2(d-1)=g-(d-2)$. (iii) The three
core neighbours have degree $3$; by (P2) they are loaded, pairwise non-adjacent,
and have $\ell=1$, i.e.\ $a=\frac16$. Hence $\tau_v=0$,
$X(v)=\frac14+3(\frac14-\frac16)=\frac12$ and $b(v)=0$.
\end{proof}

\begin{lemma}[Counting bound]\label{lem:count}
$M(G)\ge f(h):=\min_{2\le k\le h}\Bigl(\dfrac k2+\dfrac{h-k}{k-1}\Bigr)$. In
particular $h\ge4$, $M(G)\ge2$, and $M(G)\ge\frac{15}4$ whenever $h\ge10$.
\end{lemma}
\begin{proof}
This is the argument of \cite[Lemma~7]{DLS}. Order the sites
$v_1,\dots,v_h$ by non-increasing degree $d_1\ge\dots\ge d_h$. An edge at $v_i$
whose other end is not among $v_1,\dots,v_{i-1}$ has weight exactly $1/d_i$, and at
most $i-1$ edges at $v_i$ lead to earlier sites. Hence for every $k$ the edges
meeting $\{v_1,\dots,v_k\}$ have total weight at least
$\sum_{i\le k}\bigl(1-\frac{i-1}{d_i}\bigr)$. Let $k_0$ be the largest $k\le h$ with
$d_k\ge k-1$; then $k_0\ge2$, and since $d_i\ge k_0-1$ for $i\le k_0$ this weight is
at least $k_0-\frac{k_0(k_0-1)}{2(k_0-1)}=\frac{k_0}2$. For $i>k_0$ we have
$d_i\le d_{k_0+1}\le k_0-1$. By (P1), $v_i$ has a lower neighbour $z_i$; the edge
$v_iz_i$ has weight $1/d_i\ge1/(k_0-1)$, does not meet $\{v_1,\dots,v_{k_0}\}$
(their degrees are at least $k_0-1\ge d_i>d_{z_i}$), and these $h-k_0$ edges are
distinct since $v_i$ is the end of larger degree. So
$M(G)\ge\frac{k_0}2+\frac{h-k_0}{k_0-1}\ge f(h)$. The function
$f$ is non-decreasing in $h$ (each term is, and the new term $\frac{h+1}2$ for
$k=h+1$ exceeds the term $\frac h2$ for $k=h$), $f(4)=2$ and $f(10)=\frac{15}4$
(attained at $k=5$).
Finally $h=3$ is impossible: then $H=K_3$, by (P2) all three sites would be loaded
of degree $3$, and adjacent sites of degree $3$ contradict (P2).
\end{proof}

\section{The stationarity--deletion principle}\label{sec:sd}

We return to the minimal bad graph $G$ and its maximiser $p$; by
Proposition~\ref{prop:shape} and Lemma~\ref{lem:count}, $G$ is an irreducible
decorated block and $M:=M(G)\ge2$. Write $F=F(G)$, $K=K_G$ and $y=\sqrt{M-1}\ge1$.

\begin{lemma}[Stationarity]\label{lem:kkt}
$(Kp)_u=F$ for every $u\in\supp p$.
\end{lemma}
\begin{proof}
For $u,z\in V(G)$ with $p_u>0$ the vector $p+\varepsilon(\delta_z-\delta_u)$ is a
probability vector for small $\varepsilon>0$, and
$\frac{d}{d\varepsilon}F=2\bigl((Kp)_z-(Kp)_u\bigr)$ at $\varepsilon=0$. Maximality
gives $(Kp)_z\le(Kp)_u$. Hence $(Kp)_u=\max_z(Kp)_z$ for all $u\in\supp p$, and
$F=\sum_up_u(Kp)_u$ equals this common value.
\end{proof}

\begin{lemma}[Deletion inequality]\label{lem:sdel}
Let $D\subset V(G)$ be such that $G-D$ is connected with at least two vertices, and
suppose that $u\in D$ is the only vertex of $D$ that may have positive mass. Put
$t=p_u$ and $x=M(G)-M(G-D)$.
\begin{enumerate}[label=(\alph*)]
\item If $t=0$, then $x<0$.
\item If $0<t<\frac12$, then $(1-t)^2\,x<(M-1)\,t^2$.
\item In all cases, if $x\le1$ then
\[
 t\ \ge\ \theta(x^+),\qquad \theta(s):=\frac{\sqrt s}{y+\sqrt s}\quad(s\ge0),
 \qquad x^+=\max\{x,0\}.
\]
\end{enumerate}
\end{lemma}
\begin{proof}
(a) is Lemma~\ref{lem:sub}. Let $0<t<1$ and let $q$ be the probability vector
$q_z=p_z/(1-t)$ on $V(G)\setminus D$. Splitting the double sum defining
$F(G;p)$ according to whether an index equals $u$, and using
Lemma~\ref{lem:mono} for the pairs outside $D$,
\[
 F\le(1-t)^2F(G-D;q)+2t\sum_{z\ne u}p_z\,d_G(u,z)+2t^2 .
\]
By Lemma~\ref{lem:kkt}, $\sum_{z\ne u}p_zd_G(u,z)=(Kp)_u-2t=F-2t$. The graph $G-D$
has fewer vertices than $G$, hence is good: $F(G-D;q)\le M-x+1$. Therefore
\[
 (1-2t)F\le(1-t)^2(M+1-x)-2t^2 .
\]
If $t<\frac12$, inserting $F>M+1$ on the left and using
$(1-t)^2-(1-2t)=t^2$ gives $0<t^2(M+1)-2t^2-(1-t)^2x$, which is (b).
(c) If $x\le0$ there is nothing to prove. Let $0<x\le1$. Then $t>0$ by (a). If
$t<\frac12$, (b) gives $(1-t)\sqrt x<yt$, i.e.\ $t>\theta(x)$. If $t\ge\frac12$,
then $\theta(x)\le\frac12\le t$ because $\sqrt x\le1\le y$.
\end{proof}

For an irreducible decorated block put
\begin{equation}\label{eq:Theta}
\begin{gathered}
 \xi(v)=\begin{cases}\max\{X(v),b(v)\},& v\text{ loaded},\\ \max\{X(v),0\},&
 v\text{ unloaded},\end{cases}\\[4pt]
 \Theta(G)=\sum_{v\in V(H)}\theta\bigl(\xi(v)\bigr),\qquad
 \theta(s)=\frac{\sqrt s}{\sqrt{M(G)-1}+\sqrt s}.
\end{gathered}
\end{equation}
Note that $0\le\xi(v)\le1$ by Lemma~\ref{lem:X} and \eqref{eq:bleaf}.

\begin{proposition}[Mass budget]\label{prop:budget}
The minimal bad graph $G$ satisfies $\Theta(G)\le1$.
\end{proposition}
\begin{proof}
For an unloaded site $v$ apply Lemma~\ref{lem:sdel}(c) with $D=\{v\}$: then
$p_v\ge\theta(X(v)^+)=\theta(\xi(v))$. For a loaded site $v$ we have $p_v=0$
(Proposition~\ref{prop:shape}), so Lemma~\ref{lem:sdel}(c) applies to
$D=\{v'\}$, where $x=b(v)$ by \eqref{eq:bleaf}, and to $D=\{v,v'\}$, where
$x=X(v)$; both times $u=v'$. Since $\theta$ is increasing, $p_{v'}\ge\theta(\xi(v))$.
The vertices $v$
($v$ unloaded) and $v'$ ($v$ loaded) are distinct, so $\Theta(G)\le\sum_zp_z=1$.
\end{proof}

\section{The budget is exceeded}\label{sec:budget}

\begin{proposition}\label{prop:exceed}
Every irreducible decorated block $G$ satisfies $\Theta(G)>1$.
\end{proposition}

Propositions~\ref{prop:budget} and \ref{prop:exceed} contradict each other, which
proves Theorem~\ref{thm:W}. Proposition~\ref{prop:exceed} follows from the next two
lemmas together with Lemma~\ref{lem:count}: if $h\ge10$ then $M(G)\ge\frac{15}4>3.62$.

\begin{lemma}[Large blocks]\label{lem:A}
Let $G$ be an irreducible decorated block with $\Theta(G)\le1$ and put
$y=\sqrt{M(G)-1}$. Then $M(G)\le1.3+1.43\,y$, and consequently $M(G)<3.62$.
\end{lemma}

\begin{lemma}[Small blocks; computer-assisted]\label{lem:small}
Every irreducible decorated block with $h\le9$ sites satisfies
$\Theta(G)\ge1.38$.
\end{lemma}

\begin{proof}[Proof of Lemma~\ref{lem:A}]
Write $M=M(G)$; recall that $M\ge2$ by Lemma~\ref{lem:count}. Let $\mathcal P$ be
the set of loaded strict peaks of degree at least $5$ and $\Lambda$ the set of all
other loaded sites. A strict peak $z$ is the end of larger degree of each of its
$d_z-\epsilon_z$ core edges, which have weight $1/d_z$, and distinct strict peaks
share no edge; hence
\[
 \sum_{z\in\mathcal P}\frac1{d_z}\le\frac14\sum_{z\in\mathcal P}\frac{d_z-1}{d_z}\le
 \frac14M_{\rm core}.
\]
By Lemma~\ref{lem:sum} and $X(v)\le\xi(v)$,
$M_{\rm core}\le\sum_vX(v)\le\sum_v\xi(v)$. Therefore
\begin{equation}\label{eq:gain}
 M=M_{\rm core}+\sum_{v\ \rm loaded}\frac1{d_v}\ \le\ \sum_{v\in V(H)}\gamma_v,\qquad
 \gamma_v:=\frac54\,\xi(v)+\lambda_v,\quad
 \lambda_v:=\begin{cases}1/d_v,&v\in\Lambda,\\0,&\text{otherwise.}\end{cases}
\end{equation}
If $v\in\Lambda$ then $\xi(v)>0$ by Lemma~\ref{lem:upper}. Hence $\gamma_v>0$ only
if $\xi(v)>0$, and with $\rho_v:=\gamma_v/\theta(\xi(v))$ for these sites,
\[
 M\le\sum_{\xi(v)>0}\theta(\xi(v))\,\rho_v\le\Theta(G)\cdot\max_v\rho_v\le\max_v\rho_v .
\]
It remains to show $\rho_v\le1.3+1.43\,y$ for every site with $\xi(v)>0$. Put
$s=\sqrt{\xi(v)}\in(0,1]$; then
\[
 \rho_v=\Bigl(\frac54s^2+\lambda_v\Bigr)\frac{y+s}{s}
 =\frac54s^2+\lambda_v+y\Bigl(\frac54s+\frac{\lambda_v}{s}\Bigr).
\]

\emph{Case $\lambda_v=0$.} Then $\rho_v=\frac54s(y+s)\le\frac54+\frac54y$.

\emph{Case $\lambda_v=1/d$}, $d=d_v$. As a function of $s>0$, $\rho_v$ is convex,
so on an interval it is at most its larger endpoint value
$\alpha+\beta y$, where at $s^2=\sigma$
\[
 \alpha=\frac54\sigma+\frac1d,\qquad \beta=\frac54\sqrt\sigma+\frac1{d\sqrt\sigma}.
\]
By Lemma~\ref{lem:upper} the relevant values of $\sigma$ (the ends of the range of
$\xi(v)$) are the following; we list $(\alpha,\beta)$ rounded up. Recall that the
sites of $\Lambda$ of degree at least $5$ are not strict peaks.
\begin{center}\small
\begin{tabular}{lll}
\toprule
site & $\sigma$ & $(\alpha,\beta)$\\
\midrule
$d=3$ & $1/3$ & $(0.7500,\ 1.2991)$\\
$d=3$ & $2/3$ & $(1.1667,\ 1.4289)$\\
$d=4$, not a strict peak & $1/12$ & $(0.3542,\ 1.2269)$\\
$d=4$, not a strict peak & $5/6$ & $(1.2917,\ 1.4150)$\\
$d=4$, strict peak & $1/2$ & $(0.8750,\ 1.2375)$\\
$d=5$, not a strict peak & $17/20$ & $(1.2625,\ 1.3694)$\\
$d\ge5$, not a strict peak & $1/(d(d-1))$ & $(\le0.2625,\ \le1.2800)$\\
$d\ge6$, not a strict peak & $1-\frac{d-2}{d(d-1)}$ & $(\le1.2500,\ \le1.4291)$\\
\bottomrule
\end{tabular}
\end{center}
For the last two rows: at $\sigma=1/g$ with $g=d(d-1)\ge20$ we have
$\alpha=\frac5{4g}+\frac1d\le\frac1{16}+\frac15$ and
$\beta=\frac5{4\sqrt g}+\sqrt{\frac{d-1}d}<0.28+1$. At
$\sigma=\sigma_d:=1-\frac{d-2}{g}$ with $d\ge6$ we have
$\alpha=\frac54-\frac{d-6}{4g}\le\frac54$, and since $\sigma_d$ increases with $d$,
$\beta<\frac54+\frac1{6\sqrt{\sigma_6}}=\frac54+\frac1{6\sqrt{13/15}}<1.4291$.
In every case $\alpha\le1.3$ and $\beta\le1.43$.

Thus $y^2+1=M\le1.3+1.43y$, whence
$y\le\frac12\bigl(1.43+\sqrt{1.43^2+1.2}\bigr)<1.6157$ and $M<3.62$.
\end{proof}

\begin{proof}[Proof of Lemma~\ref{lem:small}]
Two programs were run on the list of all two-connected graphs $H$ on $h$ vertices,
$3\le h\le9$, generated by \texttt{nauty-geng -C} \cite{nauty}.

The first program (Appendix~\ref{app:blocks}) runs through the sets of loaded sites
that contain all core vertices of degree two (the other sets are not irreducible by
(P2)), applies the quick test (P1), which is implied by irreducibility, and then
tests irreducibility by the integer form of Lemma~\ref{lem:del}(a). For each
irreducible decoration it computes $M(G)$, $X(v)$ and $b(v)$ directly from the
definitions as integers over the common denominator $27720$, and then $\Theta(G)$ in
floating point.

The second program (Appendix~\ref{app:indep}) was written independently and uses no
formula of this paper for the test: for each of the $2^h$ sets of loaded sites it
builds $G$ explicitly and tests the definition, recomputing $M(G-e)$ for every core
edge $e$; it then computes $\Theta(G)$ in the same way and also re-checks the
statements of Lemmas~\ref{lem:struct}--\ref{lem:count} and inequality
\eqref{eq:gain} in exact arithmetic.

Both programs find the same irreducible decorations, $5\,956$ in all for $h\le9$
(counted as pairs of a core and a set of loaded sites), with the same minima; see
Table~\ref{tab:blocks}. The smallest value of $\Theta$ is $1.3845\ldots$, attained
for $H=K_2\vee\overline{K_5}$ (two adjacent vertices joined to five further
vertices) when the five vertices of degree two are loaded and at most one of the
other two is. Since $\Theta$ is a sum of at most nine terms, each computed from
exact rationals with relative error below $10^{-12}$, the bound $1.38$ is rigorous.
The numbers of cores in the second column of the table agree with the known numbers
of two-connected graphs, and in the second program every decorated block with $h$
sites arises from exactly one core and at least one of the $2^h$ sets of loaded
sites.
\end{proof}

\begin{table}[ht]
\centering\small
\begin{tabular}{rrrrrr}
\toprule
$h$ & two-connected $H$ & decorations & irreducible & $\min\Theta$ & $\min M$\\
\midrule
3 & 1 & 1 & 0 & -- & --\\
4 & 3 & 21 & 1 & 1.4065 & 2.4167\\
5 & 10 & 149 & 5 & 1.5271 & 2.7500\\
6 & 56 & 1\,873 & 12 & 1.5908 & 3.1333\\
7 & 468 & 32\,217 & 78 & 1.3846 & 3.5000\\
8 & 7\,123 & 1\,065\,605 & 494 & 1.5000 & 3.8214\\
9 & 194\,066 & 64\,360\,437 & 5\,366 & 1.6818 & 4.1250\\
\midrule
10 & 9\,743\,542 & 7\,213\,889\,689 & 81\,866 & 1.6504 & 4.4167\\

\bottomrule
\end{tabular}
\caption{Exhaustive enumeration of decorated blocks. ``Decorations'' counts the sets
of loaded sites that contain all core vertices of degree two. Rows $h\le9$ are used
in the proof; the further rows are redundant (Lemma~\ref{lem:A} covers $h\ge10$) and
serve as a check.}
\label{tab:blocks}
\end{table}

\section{Proofs of the theorems}\label{sec:proofs}

\begin{proof}[Proof of Theorem~\ref{thm:W}]
If the theorem fails, a minimal bad graph $G$ exists. It is an irreducible
decorated block (Proposition~\ref{prop:shape}) with $\Theta(G)\le1$
(Proposition~\ref{prop:budget}). If its core has at most nine sites this contradicts
Lemma~\ref{lem:small}. Otherwise $M(G)\ge\frac{15}4$ by Lemma~\ref{lem:count}, which
contradicts Lemma~\ref{lem:A}.
\end{proof}

\begin{proof}[Proof of Corollary~\ref{cor:main}]
Apply Theorem~\ref{thm:W} to the uniform measure on $\{x_1,\dots,x_k\}$ to get (b).
Part (a) is $k=2$ with $d(x_1,x_2)=\diam(G)$, and (c) is $k=n$, where
$\frac1{n^2}\sum_{x\ne y}d(x,y)=\frac{n-1}n\ad(G)$.
\end{proof}

\begin{lemma}\label{lem:RM}
Let $G$ be a graph without isolated vertices, of order $n$ and maximum degree
$\Delta$. Then
\begin{equation}\label{eq:RM}
 (\sqrt\Delta+1)\,R(G)\ \ge\ n+(\sqrt\Delta-1)\,M(G).
\end{equation}
\end{lemma}
\begin{proof}
For an edge with end degrees $a\le b$,
\[
 \Bigl(\frac1{\sqrt a}-\frac1{\sqrt b}\Bigr)^2
 =\Bigl(\sqrt{\tfrac ba}-1\Bigr)\Bigl(\frac1{\sqrt{ab}}-\frac1b\Bigr)
 \le(\sqrt\Delta-1)\Bigl(\frac1{\sqrt{ab}}-\frac1b\Bigr).
\]
Summing over the edges and using $\sum_{uv\in E}(1/d_u+1/d_v)=n$ gives
$n-2R(G)\le(\sqrt\Delta-1)(R(G)-M(G))$.
\end{proof}

\begin{lemma}[\cite{EJS}]\label{lem:path}
$\ad(G)\le\frac{n+1}3$ for every connected graph of order $n\ge2$, with equality
only for the path.
\end{lemma}
\begin{proof}
It suffices to consider a tree $T$ (pass to a spanning tree). Let $v$ be a leaf of
$T$. Listing the other vertices in breadth-first order from $v$, the $i$-th one is
at distance at most $i$, so $\sum_ud_T(u,v)\le\binom n2$, with equality only if
$T$ is a path ending at $v$. Since $W(T)=W(T-v)+\sum_ud_T(u,v)$, induction gives
$W(T)\le\binom n3+\binom n2=\binom{n+1}3$, with equality only for the path. If $G$
is not a tree, then $W(G)<W(T)$ for every spanning tree $T$ of $G$, because the ends
of an edge of $G$ outside $T$ are closer in $G$ than in $T$.
\end{proof}

\begin{proof}[Proof of Theorem~\ref{thm:F}]
Let $n\ge5$, $\mu=\ad(G)$ and $r=\sqrt{\Delta(G)}\le r_0:=\sqrt{n-1}$. By
Corollary~\ref{cor:main}(c), $M(G)\ge\bigl((n-1)\mu-(n-2)\bigr)/n$. Suppose that
$R(G)\le\mu$. Then \eqref{eq:RM} gives
\[
 (r+1)\mu\ \ge\ n+(r-1)\,\frac{(n-1)\mu-(n-2)}{n},
 \qquad\text{that is,}\qquad
 \mu\ \ge\ \frac{n^2-(r-1)(n-2)}{2n+r-1}=:\varphi(r).
\]
The function $\varphi$ is strictly decreasing in $r\ge1$ (its derivative is
$-(3n^2-4n)/(2n+r-1)^2$), so $\mu\ge\varphi(r_0)$. We claim
$\varphi(r_0)\ge\frac{n+1}3$. This is equivalent to
$n^2+2n-5\ge r_0(4n-5)$, and, after squaring, to
\[
 (n^2+2n-5)^2-(n-1)(4n-5)^2=(n-5)(n-2)(n^2-5n+5)\ \ge0,
\]
which holds for $n\ge5$. Hence $\mu\ge\varphi(r)\ge\varphi(r_0)\ge\frac{n+1}3\ge\mu$
by Lemma~\ref{lem:path}, so equality holds throughout: $n=5$, $r=r_0=2$, and $G$
is the path on five vertices, whose maximum degree is $2$, not $4$. This
contradiction proves $R(G)>\ad(G)$ for $n\ge5$.

For $n\le4$ there are nine connected graphs. Their values of $(R,\ad)$ are:
$K_2$ $(1,1)$; $P_3$ $(\sqrt2,\frac43)$; $K_3$ $(\frac32,1)$;
$P_4$ $(\frac12+\sqrt2,\frac53)$; $K_{1,3}$ $(\sqrt3,\frac32)$; $C_4$ $(2,\frac43)$;
the triangle with a pendant vertex $(\frac12+\frac1{\sqrt3}+\frac2{\sqrt6},\frac43)$;
$K_4-e$ $(\frac13+\frac4{\sqrt6},\frac76)$; $K_4$ $(2,1)$. In all cases
$R\ge\ad$, with equality only for $K_2$.
\end{proof}

\section{Remarks}\label{sec:remarks}

\begin{remark}[Sharpness of Theorem~\ref{thm:W}]\label{rem:sharp}
For a star with at least two leaves, $M=1$ and $F(G;p)=2$ for every $p$ carried by
the leaves. For a path with $n\ge3$ vertices, $M=\frac{n-1}2$ and the two-point
measure on the ends gives $F=\frac{n-1}2+1$. For $K_{2,b}$, $b\ge2$, with one pendant
vertex at each vertex of the side of size two, $M=2$, the diameter is $4$, and the two-point measure on the
pendant vertices gives $F=3$. This last graph is a decorated block that is not
irreducible (its core edges are not supportive). The quantities of \eqref{eq:Theta},
computed for it in the same way, are $\xi=1$ at the two loaded sites and $\xi=0$
elsewhere, so that $\Theta=1$: irreducibility cannot be dropped in
Proposition~\ref{prop:exceed}.
\end{remark}

\begin{remark}[The deletion inequality for all graphs]\label{rem:allgraphs}
Once Theorem~\ref{thm:W} is known, the computation in the proof of
Lemma~\ref{lem:sdel} applies to every connected graph $G$ with at least three
vertices: if $p$ is a maximiser of $F(G;\cdot)$ and $u$ is a vertex with $G-u$
connected and $t=p_u<\frac12$, then
\[
 F(G)\ \le\ M(G)+1-\frac{(1-t)^2\bigl(M(G)-M(G-u)\bigr)-(M(G)-1)\,t^2}{1-2t}.
\]
(For $t=0$ this is Lemma~\ref{lem:mono}.) Thus Theorem~\ref{thm:W} is strict as soon
as some deletable vertex carries less mass than its deletion value requires. This is
the form of the inequality that was tested numerically, see
Remark~\ref{rem:checks}(2).
\end{remark}

\begin{remark}[Independent checks]\label{rem:checks}
The following computations are not part of the proof. Programs that are not printed
in Appendix~\ref{app:code} are in the directory \texttt{extra} of the accompanying
repository.
(1) For every connected graph with at most nine vertices ($273\,192$ graphs;
$261\,080$ of order nine) the maximum $F(G)$ was computed by enumerating supports
and solving the stationarity equations, and $F(G)\le M(G)+1$ was confirmed
(Appendix~\ref{app:small}). This is legitimate because a maximiser with a minimal
number of positive coordinates has a non-singular kernel on its support. A separate C program (Appendix~\ref{app:direct}) did the same in floating point for all $11\,716\,571$ connected graphs of order ten; the $86$ graphs with $M+1-F<10^{-6}$ were then examined in exact rational arithmetic, and all of them satisfy $F=M+1$.
The inequality was also confirmed, with a heuristic maximisation that gives lower
bounds for $F$, for $1\,500$ random graphs with $8$ to $21$ vertices.
(2) The inequality of Remark~\ref{rem:allgraphs} was confirmed at the maximiser for
all connected graphs with at most eight vertices and all non-cut vertices of mass
less than $\frac12$ ($70\,315$ instances).
(3) Lemma~\ref{lem:sum}, the bound \eqref{eq:U} and the inequality \eqref{eq:gain}
were confirmed in exact arithmetic, and the bound $\rho_v\le1.3+1.43y$ in floating
point, on all $87\,822$ irreducible decorations with $h\le10$
(Appendix~\ref{app:cross}); the largest value of $\rho_v-(1.3+1.43y)$ is $-0.20$.
(4) The second program used for Lemma~\ref{lem:small} (Appendix~\ref{app:indep})
confirmed, on all $5\,956$ irreducible decorations with $h\le9$ and in exact
arithmetic, (P1)--(P3), the formula of Lemma~\ref{lem:X}, Lemma~\ref{lem:sum},
Lemma~\ref{lem:upper}, Lemma~\ref{lem:count} and \eqref{eq:gain}.
(5) All statements of Lemmas~\ref{lem:struct}--\ref{lem:count}, the inequality
\eqref{eq:gain}, the bound $\rho_v\le1.3+1.43y$ and $\Theta>1$ were confirmed from
the definitions, with rational arithmetic, on 5\,640 further irreducible decorated
blocks with up to 111 sites, found by random search among cores consisting of a
few vertices of large degree joined through many vertices of degree two or three,
and among complete multipartite cores (Appendix~\ref{app:big}).
(6) A separate AI review pass re-derived Sections~\ref{sec:prelim}--\ref{sec:proofs}
and tested each lemma numerically with its own programs; the second program of Lemma~\ref{lem:small} and the check for order ten in (1) come from that pass. It is
not a substitute for refereeing.
\end{remark}

\begin{remark}[The diameter theorem]
Corollary~\ref{cor:main}(a) is the theorem of \cite{DLS}. The proof given here is
independent of the proof in \cite{DLS}, from which only the padding idea and the
counting argument of Lemma~\ref{lem:count} are borrowed.
\end{remark}

\begin{remark}[Erratum to an earlier manuscript]
An earlier manuscript of this series (\emph{Breadth-first routing trees and
Fajtlowicz's average-distance conjecture}, 15 November 2026) proposed a
``deepest-vertex inequality'' as Conjecture~3.4 and verified it for all graphs with
at most nine vertices. As stated, for \emph{every} vertex at maximum distance from
the root, that conjecture is false: in $K_2\vee C_{41}$ ($43$ vertices), rooted at
one of the two vertices of degree $42$, the other one is a deepest vertex and the
inequality reads $1.906\ldots\ge3.952\ldots$. The breadth-first routing inequality
proposed there (its Conjecture~1.2) is not affected by this example and remains
open; it is not implied by Theorem~\ref{thm:F}.
\end{remark}

\subsection{Audit against six acceptance criteria}\label{sec:audit}
The work was carried out against six criteria fixed in advance (G1--G6 below). We
record how the result stands with respect to each of them.
\begin{center}\small
\begin{tabular}{p{0.27\linewidth}p{0.66\linewidth}}
\toprule
Criterion & Assessment\\
\midrule
G1 a named problem that existed before this work & Fajtlowicz's average-distance conjecture \cite{Faj} (1988), in the form quoted in \cite{CPS,YL,LiShi}.\\
G2 the problem is not changed & Theorem~\ref{thm:F} is the original statement, for all connected graphs and with the original quantities $R$ and $\ad$. The weighted inequality is an intermediate strengthening, not a replacement.\\
G3 the problem was open & In the literature we could find \cite{CPS,YL,LiShi,Dalfo,VMRC189}, and in searches of the arXiv listings on 20 November 2026, only special classes were settled. A search cannot exclude a proof that is not indexed.\\
G4 unconditional and complete & A complete proof is given, with no hypothesis and no restriction on the graph. One lemma (Lemma~\ref{lem:small}) is a finite exhaustive computation whose source is printed in Appendix~\ref{app:code}. The manuscript is unrefereed: independent verification is required before the result is relied upon.\\
G5 not a restatement & The stationarity--deletion inequality and the resulting mass budget are new estimates; they are what closes the case of unbounded degree and diameter that was left open in \cite{VMRC189}.\\
G6 no earlier identical solution & No proof of the conjecture and no use of the stationarity--deletion argument was found. Priority cannot be certified by the absence of search results.\\
\bottomrule
\end{tabular}
\end{center}

\begin{thebibliography}{99}
\bibitem{CPS} M.~Cygan, M.~Pilipczuk and R.~\v{S}krekovski, Relation between Randi\'c index and average distance of trees, \emph{MATCH Commun.\ Math.\ Comput.\ Chem.} \textbf{66} (2011), 605--612.
\bibitem{Dalfo} C.~Dalf\'o, On the Randi\'c index of graphs, \emph{Discrete Math.} \textbf{342} (2019), 2792--2796.
\bibitem{DLS} Z.~Dvo\v{r}\'ak, B.~Lidick\'y and R.~\v{S}krekovski, Randi\'c index and the diameter of a graph, \emph{European J.\ Combin.} \textbf{32} (2011), 434--442; preprint KAM--DIMATIA Series 2010--496.
\bibitem{EJS} R.~C.~Entringer, D.~E.~Jackson and D.~A.~Snyder, Distance in graphs, \emph{Czechoslovak Math.\ J.} \textbf{26} (1976), 283--296.
\bibitem{Faj} S.~Fajtlowicz, On conjectures of Graffiti, \emph{Discrete Math.} \textbf{72} (1988), 113--118.
\bibitem{LiShi} X.~Li and Y.~Shi, Randi\'c index, diameter and the average distance, arXiv:0906.5230 (2009).
\bibitem{nauty} B.~D.~McKay and A.~Piperno, Practical graph isomorphism, II, \emph{J.\ Symbolic Comput.} \textbf{60} (2014), 94--112.
\bibitem{Randic} M.~Randi\'c, On characterization of molecular branching, \emph{J.\ Amer.\ Chem.\ Soc.} \textbf{97} (1975), 6609--6615.
\bibitem{VMRC189} Keio and Shaitetsu, \emph{Collected research papers: the Randi\'c index, graph distances, and rational Lambert series}, collected edition, November 2026, papers 06--14; item [189] of the Verifiable Mathematical Reasoning Corpus, \url{https://github.com/QEEDDD/Verifiable-Mathematical-Reasoning-Corpus}. Unrefereed.
\bibitem{YL} Z.~You and B.~Liu, On a conjecture between Randi\'c index and average distance of unicyclic graphs, \emph{Filomat} \textbf{28} (2014), 767--773.
\end{thebibliography}

\appendix
\section{Programs}\label{app:code}
Graphs are generated with \texttt{nauty} 2.8.8. Commands:
\begin{lstlisting}
gcc -O3 -o irreducible_blocks irreducible_blocks.c -lm
for h in 3 4 5 6 7 8 9 10; do nauty-geng -C -q $h | ./irreducible_blocks; done
gcc -O3 -o lemma_crosscheck lemma_crosscheck.c -lm
for h in 4 5 6 7 8 9 10; do nauty-geng -C -q $h | ./lemma_crosscheck; done
gcc -O3 -o independent_enum independent_enum.c -lm
for h in 3 4 5 6 7 8 9; do nauty-geng -C -q $h | ./independent_enum; done
for n in 2 3 4 5 6 7 8 9; do nauty-geng -c -q $n | python3 wc_small.py; done
gcc -O3 -o direct_check direct_check.c -lm
for k in $(seq 0 15); do nauty-geng -c -q 10 $k/16 | ./direct_check 2>>tight_n10.g6; done
python3 extra/exact_tight.py < tight_n10.g6
for s in 1 2 3 4; do python3 bigblocks.py $s 3000; done
\end{lstlisting}
(For the larger cases the input was split with the \texttt{res/mod} argument of
\texttt{geng}, e.g.\ \texttt{nauty-geng -C -q 9 0/4}.)

\subsection{Enumeration of irreducible decorated blocks (Lemma~\ref{lem:small})}\label{app:blocks}
\begin{lstlisting}[language=C]
/* Enumerate decorated two-connected blocks G = H + leaves (at most one leaf per site),
   H read as graph6 from stdin (output of: nauty-geng -C -q h).
   A decoration is IRREDUCIBLE if every core edge uv is supportive:  M(G-uv) > M(G),
   M(G) = sum over edges 1/max(deg,deg).  For every irreducible decoration compute exactly
   (integers over DEN = 27720) the deletion values
       X_v = M(G) - M(G - v - leaf(v))   (every site v),
       b_v = M(G) - M(G - leaf(v))       (loaded sites),
   xi_v = max(X_v, b_v, 0),  and  Theta = sum_v theta(xi_v),
   theta(x) = sqrt(x)/(sqrt(M-1)+sqrt(x)).   Reports the minimum of Theta.            */
#include <stdio.h>
#include <stdlib.h>
#include <string.h>
#include <math.h>
#define MAXH 12
#define DEN 27720L   /* lcm(1,...,12): every degree that can occur for h<=11 divides DEN */
static int h; static unsigned adj[MAXH]; static int kdeg[MAXH];
static int eu[64], ev[64], ne;
static void parse(const char *s){
  h = s[0]-63; int bit=5; const char *p=s+1; int val=*p-63;
  for(int i=0;i<h;i++) adj[i]=0;
  for(int j=1;j<h;j++) for(int i=0;i<j;i++){
    if(bit<0){p++; val=*p-63; bit=5;}
    if((val>>bit)&1){adj[i]|=1u<<j; adj[j]|=1u<<i;}
    bit--; }
  ne=0; for(int i=0;i<h;i++){ kdeg[i]=__builtin_popcount(adj[i]);
    for(int j=i+1;j<h;j++) if(adj[i]>>j&1){eu[ne]=i; ev[ne]=j; ne++;} }
}
/* M (times DEN) of the graph: sites in `alive`, core edges among alive sites, leaves at sites in lm&alive */
static long Mof(unsigned alive, unsigned lm){
  int d[MAXH]; long s=0;
  for(int v=0;v<h;v++) d[v] = (alive>>v&1) ? __builtin_popcount(adj[v]&alive) + ((lm>>v)&1) : 0;
  for(int e=0;e<ne;e++){ int u=eu[e], v=ev[e];
    if((alive>>u&1) && (alive>>v&1)){ int m = d[u]>d[v]?d[u]:d[v]; s += DEN/m; } }
  for(int v=0;v<h;v++) if((alive>>v&1) && (lm>>v&1)) s += DEN/d[v];
  return s;
}
int main(void){
  char line[128]; long cores=0, decos=0, irr=0, bad=0; double minT=1e9, minM=1e9; char arg[128]=""; unsigned argm=0;
  long cnt_p1=0;
  while(fgets(line,sizeof line,stdin)){
    line[strcspn(line,"\n")]=0; parse(line); cores++;
    unsigned full=(1u<<h)-1, forced=0;
    for(int v=0;v<h;v++) if(kdeg[v]==2) forced|=1u<<v;
    unsigned freem = full & ~forced;
    /* iterate over all subsets sub of freem */
    unsigned sub=0;
    do{
      unsigned lm = forced|sub;
      {
        decos++;
        int d[MAXH]; for(int v=0;v<h;v++) d[v]=kdeg[v]+((lm>>v)&1);
        /* P1: every site has a strictly lower neighbour */
        int ok=1; int ell[MAXH];
        for(int v=0;v<h && ok;v++){
          int l=(lm>>v)&1; unsigned a=adj[v];
          while(a){int w=__builtin_ctz(a); a&=a-1; if(d[w]<d[v]) l++;}
          ell[v]=l; if(l==0) ok=0; }
        if(ok){ cnt_p1++;
          for(int e=0;e<ne && ok;e++){ int u=eu[e], v=ev[e];
            long bu = ell[u] - (d[v]<d[u]), bv = ell[v] - (d[u]<d[v]);
            long gu=(long)d[u]*(d[u]-1), gv=(long)d[v]*(d[v]-1); long mx = d[u]>d[v]?d[u]:d[v];
            /* supportive iff bu/gu + bv/gv > 1/mx */
            if(!((bu*gv + bv*gu)*mx > gu*gv)) ok=0; }
        }
        if(ok){ irr++;
          long M = Mof(full,lm);
          double Mf=(double)M/DEN, y=sqrt(Mf-1.0), T=0;
          for(int v=0;v<h;v++){
            long X = M - Mof(full & ~(1u<<v), lm & ~(1u<<v));
            long xi = X>0?X:0;
            if(lm>>v&1){ long b = M - Mof(full, lm & ~(1u<<v)); if(b>xi) xi=b; }
            if(xi>0){ double s=sqrt((double)xi/DEN); T += s/(y+s); }
          }
          if(T<minT){minT=T; strcpy(arg,line); argm=lm;}
          if(Mf<minM) minM=Mf;
          if(T<=1.0+1e-9){ bad++; printf("THETA<=1 %s mask=%u M=%.6f Theta=%.6f\n",line,lm,Mf,T); }
        }
      }
      sub = (sub - freem) & freem;   /* next subset of freem */
    } while(sub);
  }
  printf("h=%d cores=%ld decorations=%ld passP1=%ld irreducible=%ld minTheta=%.6f (at %s mask %u) minM=%.6f Theta<=1:%ld\n",
         h,cores,decos,cnt_p1,irr,minT,arg,argm,minM,bad);
  return 0;
}
\end{lstlisting}
Output for $h\le9$:
\begin{lstlisting}
h=3 cores=1 decorations=1 passP1=1 irreducible=0 minTheta=1000000000.000000 (at  mask 0) minM=1000000000.000000 Theta<=1:0
h=4 cores=3 decorations=21 passP1=3 irreducible=1 minTheta=1.406515 (at C^ mask 15) minM=2.416667 Theta<=1:0
h=5 cores=10 decorations=149 passP1=17 irreducible=5 minTheta=1.527081 (at DF{ mask 31) minM=2.750000 Theta<=1:0
h=6 cores=56 decorations=1873 passP1=168 irreducible=12 minTheta=1.590805 (at E?~w mask 15) minM=3.133333 Theta<=1:0
h=7 cores=468 decorations=32217 passP1=2283 irreducible=78 minTheta=1.384595 (at F?B~w mask 31) minM=3.500000 Theta<=1:0
h=8 cores=7123 decorations=1065605 passP1=57962 irreducible=494 minTheta=1.500000 (at G??F~w mask 63) minM=3.821429 Theta<=1:0
h=9 cores=194066 decorations=64360437 passP1=2689544 irreducible=5366 minTheta=1.681772 (at H???F~} mask 127) minM=4.125000 Theta<=1:0
\end{lstlisting}
(For $h=3$ there is no irreducible decoration, and the two minima printed are the
initial sentinel values.) Output for $h=10$ (input split into the four parts \texttt{0/4}, \dots, \texttt{3/4};
the counts add up to the row of Table~\ref{tab:blocks}):
\begin{lstlisting}
h=10 cores=2397411 decorations=1718983912 passP1=58428451 irreducible=23708 minTheta=1.650367 (at I???FreFw mask 639) minM=4.416667 Theta<=1:0
h=10 cores=2216372 decorations=1673411265 passP1=54354154 irreducible=12337 minTheta=1.812079 (at I??CFF\nW mask 383) minM=4.583333 Theta<=1:0
h=10 cores=2525186 decorations=1906671944 passP1=62909060 irreducible=14295 minTheta=1.842464 (at I??EEC~Nw mask 127) minM=4.533333 Theta<=1:0
h=10 cores=2604573 decorations=1914822568 passP1=62783569 irreducible=31526 minTheta=1.894820 (at I?AAFB|}w mask 511) minM=4.660714 Theta<=1:0
\end{lstlisting}


\subsection{Independent re-implementation (Lemma~\ref{lem:small})}\label{app:indep}
\begin{lstlisting}[language=C]
/* Independent re-implementation of the enumeration of irreducible decorated blocks
   (written separately from irreducible_blocks.c, from the definitions only).
   Reads graph6 (two-connected cores from geng -C).  For EVERY subset of loaded sites
   (no forcing, no P1 pre-test) builds G explicitly as an adjacency matrix on up to 2h vertices,
   and tests the definition: every core edge e satisfies M(G-e) > M(G), M recomputed from scratch.
   All quantities exact: integers over DEN = lcm(1..12) = 27720.
   For irreducible blocks: X(v), b(v) from the definition, xi, Theta (long double),
   plus checks: Lemma 4.3, (4.2), Lemma 4.4(i)-(iii), Lemma 4.5, (6.1), rho bound.            */
#include <stdio.h>
#include <string.h>
#include <math.h>
#define MAXN 24
#define DEN 27720LL
static int n;                 /* number of vertices of G currently */
static int A[MAXN][MAXN];
static int alive[MAXN];
static long long Mcalc(void){
  int deg[MAXN]; long long s=0;
  for(int i=0;i<n;i++){ deg[i]=0; if(!alive[i]) continue; for(int j=0;j<n;j++) if(alive[j]&&A[i][j]) deg[i]++; }
  for(int i=0;i<n;i++) if(alive[i]) for(int j=i+1;j<n;j++) if(alive[j]&&A[i][j]){
     int m=deg[i]>deg[j]?deg[i]:deg[j]; s+=DEN/m; }
  return s;
}
int main(void){
  char line[256]; long long cores=0, decos=0, irr=0, viol=0, viol2=0; long double minT=1e9L; double minM=1e9; long long thle1=0;
  double worst=-1e9; char argline[256]=""; unsigned argmask=0; int h=0; long long mintie=0;
  while(fgets(line,sizeof line,stdin)){
    line[strcspn(line,"\n")]=0; if(!line[0]) continue; cores++;
    h=line[0]-63; int H[12][12]; memset(H,0,sizeof H);
    { int k=0; /* bit index */
      for(int j=1;j<h;j++) for(int i=0;i<j;i++){
        int byte=line[1+k/6]-63; int bit=5-(k%6); if((byte>>bit)&1){H[i][j]=H[j][i]=1;} k++; } }
    int ce_u[64], ce_v[64], nce=0;
    for(int i=0;i<h;i++) for(int j=i+1;j<h;j++) if(H[i][j]){ce_u[nce]=i; ce_v[nce]=j; nce++;}
    for(unsigned mask=0; mask<(1u<<h); mask++){
      decos++;
      /* build G: sites 0..h-1, leaf of site v is vertex h+v (alive only if loaded) */
      n=2*h; memset(A,0,sizeof A);
      for(int i=0;i<h;i++) for(int j=0;j<h;j++) A[i][j]=H[i][j];
      for(int v=0;v<h;v++){ alive[v]=1; alive[h+v]=(mask>>v)&1; if(alive[h+v]){A[v][h+v]=A[h+v][v]=1;} }
      long long M=Mcalc(); int ok=1;
      for(int e=0;e<nce&&ok;e++){ int u=ce_u[e], v=ce_v[e];
        A[u][v]=A[v][u]=0; long long M2=Mcalc(); A[u][v]=A[v][u]=1;
        if(!(M2>M)) ok=0; }
      if(!ok) continue;
      irr++;
      int d[12]; for(int v=0;v<h;v++){ d[v]=(mask>>v)&1; for(int w=0;w<h;w++) d[v]+=H[v][w]; }
      long long Mcore=0; for(int e=0;e<nce;e++){int m=d[ce_u[e]]>d[ce_v[e]]?d[ce_u[e]]:d[ce_v[e]]; Mcore+=DEN/m;}
      double Mf=(double)M/DEN; long double y=sqrtl((long double)M/DEN-1.0L), T=0;
      long long sumX=0, gam4=0; /* gam4 = 4*DEN*sum gamma */
      long long P5sum=0;
      for(int v=0;v<h;v++){
        int loaded=(mask>>v)&1;
        alive[v]=0; int sv=alive[h+v]; alive[h+v]=0;
        long long X=M-Mcalc(); alive[v]=1; alive[h+v]=sv;
        long long xi=X>0?X:0, b=0;
        int nl=0,neq=0,nh=0; for(int w=0;w<h;w++) if(H[v][w]){ if(d[w]<d[v]) nl++; else if(d[w]==d[v]) neq++; else nh++; }
        int peak=(neq==0&&nh==0);
        if(X>DEN) viol++;
        if(nl+loaded==0) viol++;                 /* P1 */
        if(d[v]<3) viol++;                       /* P2 */
        long long lam4=0;
        if(loaded){
          alive[h+v]=0; b=M-Mcalc(); alive[h+v]=1;
          if(b<0) viol++;
          long long g=(long long)d[v]*(d[v]-1);
          if(b*g != (long long)(d[v]-nl-1)*DEN) viol++;         /* (4.1): b = (d-ell)/g */
          if(X*g > ((long long)(nl+1)*(d[v]-1)+nh)*DEN) viol++; /* (4.2) */
          if(b>xi) xi=b;
          if(d[v]==3){ if(!(3*xi>=DEN && 3*xi<=2*DEN)) viol++; if(nl||neq) viol++; }
          else if(!peak){ if(!(xi*g>=DEN && xi*g<=(g-(d[v]-2))*DEN)) viol++; }
          else if(d[v]==4){ if(2*xi!=DEN) viol++; }
          if(peak && d[v]>=5) P5sum+=DEN/d[v]; else lam4=4*DEN/d[v];
        }
        sumX+=X; gam4+=5*xi+lam4;
        if(xi>0){ long double s=sqrtl((long double)xi/DEN); long double th=s/(y+s); T+=th;
          double rho=(double)(((long double)(5*xi+lam4)/(4.0L*DEN))/th); double bd=1.3+1.43*(double)y;
          if(rho-bd>worst) worst=rho-bd; }
        else if(lam4>0) viol++;
      }

      { /* extra exact checks: (P2), (P3), Lemma 4.2 */
        int ell[12]; long long g_[12], a_[12], b_[12];
        for(int v=0;v<h;v++){ ell[v]=(mask>>v)&1; for(int w=0;w<h;w++) if(H[v][w]&&d[w]<d[v]) ell[v]++;
          g_[v]=(long long)d[v]*(d[v]-1); a_[v]=DEN/g_[v]*ell[v]; b_[v]=DEN/g_[v]*(d[v]-ell[v]); if(DEN%g_[v]) viol2++; }
        for(int v=0;v<h;v++) if(d[v]==3){ if(!((mask>>v)&1)) viol2++; for(int w=0;w<h;w++) if(H[v][w]&&d[w]<4) viol2++; }
        for(int e=0;e<nce;e++){ int u=ce_u[e], v=ce_v[e]; if(d[u]>d[v]){int t=u;u=v;v=t;}
          if(d[u]<d[v]){ if(!(a_[u]>b_[v])) viol2++; } else { if(!(a_[u]>=b_[v] && a_[v]>=b_[u])) viol2++; } }
        for(int v=0;v<h;v++){
          long long form=((mask>>v)&1)?DEN/d[v]:0;
          for(int w=0;w<h;w++) if(H[v][w]){ if(d[w]>d[v]) form+=b_[w]; else if(d[w]==d[v]) form+=b_[w]-DEN/g_[w]; else form+=DEN/d[v]-a_[w]; }
          for(int w=0;w<h;w++) if(H[v][w]) for(int z=w+1;z<h;z++) if(H[v][z]&&H[w][z]&&d[w]==d[z]) form-=DEN/g_[w];
          alive[v]=0; int sv=alive[h+v]; alive[h+v]=0; long long X=M-Mcalc(); alive[v]=1; alive[h+v]=sv;
          if(form!=X) viol2++; }
      }
      if(sumX<Mcore) viol++;                    /* Lemma 4.3 */
      if(4*P5sum>Mcore) viol++;
      if(4*M>gam4) viol++;                      /* (6.1) */
      { /* Lemma 4.5: M >= f(h) */
        int okf=0; long long best=-1;
        for(int t=2;t<=h;t++){ long long val=DEN*t/2 + DEN*(h-t)/(t-1); /* exact since DEN divisible by 2 and by t-1<=11 */
          if(best<0||val<best) best=val; }
        if(M>=best) okf=1; if(!okf) viol++; if(M<2*DEN) viol++; }
      if(T<minT-1e-12L){minT=T; strcpy(argline,line); argmask=mask; mintie=1;} else if(fabsl(T-minT)<=1e-12L) mintie++;
      if(Mf<minM) minM=Mf;
      if(T<=1.0L+1e-9L) thle1++;
    }
  }
  printf("h=%d cores=%lld all-decorations=%lld irreducible=%lld minTheta=%.6Lf (at %s mask %u; ties %lld) minM=%.6f Theta<=1: %lld lemma-violations=%lld P2/P3/Lemma4.2-violations=%lld max(rho-bound)=%.4f\n",
     h,cores,decos,irr,minT,argline,argmask,mintie,minM,thle1,viol,viol2,worst);
  return 0;
}
\end{lstlisting}
Output for $h\le8$, and for $h=9$ in four parts (the counts add up to $194\,066$
cores and $5\,366$ irreducible decorations):
\begin{lstlisting}
h=3 cores=1 all-decorations=8 irreducible=0 minTheta=1000000000.000000 (at  mask 0; ties 0) minM=1000000000.000000 Theta<=1: 0 lemma-violations=0 P2/P3/Lemma4.2-violations=0 max(rho-bound)=-1000000000.0000
h=4 cores=3 all-decorations=48 irreducible=1 minTheta=1.406515 (at C^ mask 15; ties 1) minM=2.416667 Theta<=1: 0 lemma-violations=0 P2/P3/Lemma4.2-violations=0 max(rho-bound)=-0.5729
h=5 cores=10 all-decorations=320 irreducible=5 minTheta=1.527081 (at DF{ mask 31; ties 1) minM=2.750000 Theta<=1: 0 lemma-violations=0 P2/P3/Lemma4.2-violations=0 max(rho-bound)=-0.5870
h=6 cores=56 all-decorations=3584 irreducible=12 minTheta=1.590805 (at E?~w mask 15; ties 3) minM=3.133333 Theta<=1: 0 lemma-violations=0 P2/P3/Lemma4.2-violations=0 max(rho-bound)=-0.4598
h=7 cores=468 all-decorations=59904 irreducible=78 minTheta=1.384595 (at F?B~w mask 31; ties 3) minM=3.500000 Theta<=1: 0 lemma-violations=0 P2/P3/Lemma4.2-violations=0 max(rho-bound)=-0.2028
h=8 cores=7123 all-decorations=1823488 irreducible=494 minTheta=1.500000 (at G??F~w mask 63; ties 4) minM=3.821429 Theta<=1: 0 lemma-violations=0 P2/P3/Lemma4.2-violations=0 max(rho-bound)=-0.2040
h=9 cores=47669 all-decorations=24406528 irreducible=1327 minTheta=1.681772 (at H???F~} mask 127; ties 4) minM=4.125000 Theta<=1: 0 lemma-violations=0 P2/P3/Lemma4.2-violations=0 max(rho-bound)=-0.2055
h=9 cores=35374 all-decorations=18111488 irreducible=683 minTheta=1.849527 (at H?AELvz mask 511; ties 1) minM=4.303571 Theta<=1: 0 lemma-violations=0 P2/P3/Lemma4.2-violations=0 max(rho-bound)=-0.2058
h=9 cores=70139 all-decorations=35911168 irreducible=1849 minTheta=1.872558 (at H?B@c\| mask 127; ties 2) minM=4.316667 Theta<=1: 0 lemma-violations=0 P2/P3/Lemma4.2-violations=0 max(rho-bound)=-0.2054
h=9 cores=40884 all-decorations=20932608 irreducible=1507 minTheta=1.933889 (at H?bDLj^ mask 127; ties 2) minM=4.392857 Theta<=1: 0 lemma-violations=0 P2/P3/Lemma4.2-violations=0 max(rho-bound)=-0.2055
\end{lstlisting}

\subsection{Cross-check of the analytic lemma}\label{app:cross}
\begin{lstlisting}[language=C]
/* Cross-check of the analytic lemma on every irreducible decorated block (same enumeration as irreducible_blocks.c).
   Exact integer checks (denominator 27720):
     (I1) sum_v X_v >= M_core
     (I2) X_v <= U_v = ((n_l+1)(d-1)+n_h)/(d(d-1))  for loaded non-peak sites
     (I3) M <= sum_v gamma_v,  gamma_v = (5/4) xi_v + lambda_v
   Floating check:
     (I4) rho_v = gamma_v (y+sqrt(xi_v))/sqrt(xi_v) <= 1.3 + 1.43 y  for every site with xi_v>0 */
#include <stdio.h>
#include <stdlib.h>
#include <string.h>
#include <math.h>
#define MAXH 12
#define DEN 27720L
static int h; static unsigned adj[MAXH]; static int kdeg[MAXH]; static int eu[64], ev[64], ne;
static void parse(const char *s){ h=s[0]-63; int bit=5; const char*p=s+1; int val=*p-63;
  for(int i=0;i<h;i++) adj[i]=0;
  for(int j=1;j<h;j++) for(int i=0;i<j;i++){ if(bit<0){p++; val=*p-63; bit=5;} if((val>>bit)&1){adj[i]|=1u<<j; adj[j]|=1u<<i;} bit--; }
  ne=0; for(int i=0;i<h;i++){ kdeg[i]=__builtin_popcount(adj[i]); for(int j=i+1;j<h;j++) if(adj[i]>>j&1){eu[ne]=i; ev[ne]=j; ne++;} } }
static long Mof(unsigned alive, unsigned lm, long *core){ int d[MAXH]; long s=0, c=0;
  for(int v=0;v<h;v++) d[v]=(alive>>v&1)?__builtin_popcount(adj[v]&alive)+((lm>>v)&1):0;
  for(int e=0;e<ne;e++){int u=eu[e],v=ev[e]; if((alive>>u&1)&&(alive>>v&1)){int m=d[u]>d[v]?d[u]:d[v]; c+=DEN/m;}}
  s=c; for(int v=0;v<h;v++) if((alive>>v&1)&&(lm>>v&1)) s+=DEN/d[v];
  if(core) *core=c; return s; }
int main(void){ char line[128]; long irr=0, f1=0,f2=0,f3=0,f4=0; double worst=-1e9, maxratio=0;
  while(fgets(line,sizeof line,stdin)){ line[strcspn(line,"\n")]=0; parse(line);
    unsigned full=(1u<<h)-1, forced=0; for(int v=0;v<h;v++) if(kdeg[v]==2) forced|=1u<<v;
    unsigned freem=full&~forced, sub=0;
    do{ unsigned lm=forced|sub;
      { int d[MAXH], ell[MAXH], ok=1; for(int v=0;v<h;v++) d[v]=kdeg[v]+((lm>>v)&1);
        for(int v=0;v<h&&ok;v++){ int l=(lm>>v)&1; unsigned a=adj[v]; while(a){int w=__builtin_ctz(a); a&=a-1; if(d[w]<d[v]) l++;} ell[v]=l; if(!l) ok=0; }
        for(int e=0;e<ne&&ok;e++){ int u=eu[e],v=ev[e]; long bu=ell[u]-(d[v]<d[u]), bv=ell[v]-(d[u]<d[v]);
          long gu=(long)d[u]*(d[u]-1), gv=(long)d[v]*(d[v]-1), mx=d[u]>d[v]?d[u]:d[v]; if(!((bu*gv+bv*gu)*mx>gu*gv)) ok=0; }
        if(ok){ irr++; long Mc; long M=Mof(full,lm,&Mc); double Mf=(double)M/DEN, y=sqrt(Mf-1.0);
          long sumX=0; long gam4=0; /* 4*DEN*sum gamma */
          for(int v=0;v<h;v++){
            long X=M-Mof(full&~(1u<<v), lm&~(1u<<v), NULL); sumX+=X;
            long xi=X>0?X:0; int loaded=(lm>>v)&1; long b=0;
            int nl=0,nh=0,ne_=0; unsigned a=adj[v]; while(a){int w=__builtin_ctz(a); a&=a-1; if(d[w]<d[v]) nl++; else if(d[w]>d[v]) nh++; else ne_++;}
            int peak=(nh==0&&ne_==0);
            if(loaded){ b=M-Mof(full, lm&~(1u<<v), NULL); if(b>xi) xi=b;
              if(!peak){ /* (I2): X*g <= ((nl+1)(d-1)+nh)*DEN */
                long g=(long)d[v]*(d[v]-1); if(X*g > ((long)(nl+1)*(d[v]-1)+nh)*DEN) f2++; } }
            long lam4 = (loaded && !(peak && d[v]>=5)) ? 4*DEN/d[v] : 0;   /* 4*DEN*lambda ; DEN/d integer */
            long g4 = 5*xi + lam4;  gam4 += g4;
            if(xi>0){ double s=sqrt((double)xi/DEN); double rho=((double)g4/(4.0*DEN))*(y+s)/s; double bound=1.3+1.43*y;
              if(rho>bound+1e-9) f4++; if(rho-bound>worst) worst=rho-bound; if(rho/(Mf)>maxratio) maxratio=rho/Mf; }
            else if(lam4>0) f3+=1000000; /* lambda>0 must imply xi>0 */
          }
          if(sumX<Mc) f1++;
          if(4*M>gam4) f3++;
        } }
      sub=(sub-freem)&freem; }while(sub); }
  printf("h=%d irreducible=%ld  violations: I1=%ld I2=%ld I3=%ld I4=%ld   max(rho-(1.3+1.43y))=%.4f\n",h,irr,f1,f2,f3,f4,worst);
  return 0; }
\end{lstlisting}
Output for $h\le9$, and for $h=10$ in two parts:
\begin{lstlisting}
h=4 irreducible=1  violations: I1=0 I2=0 I3=0 I4=0   max(rho-(1.3+1.43y))=-0.5729
h=5 irreducible=5  violations: I1=0 I2=0 I3=0 I4=0   max(rho-(1.3+1.43y))=-0.5870
h=6 irreducible=12  violations: I1=0 I2=0 I3=0 I4=0   max(rho-(1.3+1.43y))=-0.4598
h=7 irreducible=78  violations: I1=0 I2=0 I3=0 I4=0   max(rho-(1.3+1.43y))=-0.2028
h=8 irreducible=494  violations: I1=0 I2=0 I3=0 I4=0   max(rho-(1.3+1.43y))=-0.2040
h=9 irreducible=5366  violations: I1=0 I2=0 I3=0 I4=0   max(rho-(1.3+1.43y))=-0.2054
h=10 irreducible=38003  violations: I1=0 I2=0 I3=0 I4=0   max(rho-(1.3+1.43y))=-0.2067
h=10 irreducible=43863  violations: I1=0 I2=0 I3=0 I4=0   max(rho-(1.3+1.43y))=-0.2069
\end{lstlisting}

\subsection{Direct small-order check of Theorems~\ref{thm:W} and \ref{thm:F}}\label{app:small}
\begin{lstlisting}[language=Python]
#!/usr/bin/env python3
# Independent small-order check of  F(G) <= M(G)+1  and of  R(G) >= mu(G).
# Usage:  nauty-geng -c -q n | python3 wc_small.py
# F(G) = max over the simplex of p^T K p, K = distance matrix with diagonal 2, is computed
# by enumerating supports S and solving K_S w = 1 (every maximiser is of this form on its support).
import sys, itertools, math
import numpy as np, networkx as nx
def Fmax(G):
    nodes = list(G.nodes()); n = len(nodes); idx = {v: i for i, v in enumerate(nodes)}
    K = np.zeros((n, n))
    for s, dist in nx.all_pairs_shortest_path_length(G):
        for t, dd in dist.items(): K[idx[s], idx[t]] = dd
    D = K.copy(); np.fill_diagonal(K, 2.0)
    best = 2.0
    for r in range(2, n + 1):
        for S in itertools.combinations(range(n), r):
            KS = K[np.ix_(S, S)]
            try: w = np.linalg.solve(KS, np.ones(r))
            except np.linalg.LinAlgError: continue
            tot = w.sum()
            if abs(tot) < 1e-13: continue
            p = w / tot
            if (p < -1e-10).any(): continue
            best = max(best, float(p @ KS @ p))
    return best, D
cnt = 0; minslack = 1e9; minRmu = 1e9; fail = 0
for line in sys.stdin:
    G = nx.from_graph6_bytes(line.strip().encode()); n = G.number_of_nodes(); cnt += 1
    d = dict(G.degree())
    M = sum(1.0 / max(d[u], d[v]) for u, v in G.edges())
    R = sum(1.0 / math.sqrt(d[u] * d[v]) for u, v in G.edges())
    F, D = Fmax(G); mu = D.sum() / (n * (n - 1))
    minslack = min(minslack, M + 1 - F); minRmu = min(minRmu, R - mu)
    if F > M + 1 + 1e-9 or R < mu - 1e-12: fail += 1; print("FAIL", line.strip())
print(f"graphs={cnt} failures={fail} min(M+1-F)={minslack:.6f} min(R-mu)={minRmu:.6f}")
\end{lstlisting}
Output (a failure is reported if $F>M+1+10^{-9}$ or $R<\operatorname{ad}-10^{-12}$;
the value $-0.000000$ is a rounding error at graphs with $F=M+1$; the prefixes
\texttt{n=}\dots\ were added by the calling loop):
\begin{lstlisting}
n=2 graphs=1 failures=0 min(M+1-F)=0.000000 min(R-mu)=0.000000
n=3 graphs=2 failures=0 min(M+1-F)=0.000000 min(R-mu)=0.080880
n=4 graphs=6 failures=0 min(M+1-F)=0.000000 min(R-mu)=0.232051
n=5 graphs=21 failures=0 min(M+1-F)=0.000000 min(R-mu)=0.400000
n=6 graphs=112 failures=0 min(M+1-F)=0.000000 min(R-mu)=0.569401
n=7 graphs=853 failures=0 min(M+1-F)=-0.000000 min(R-mu)=0.735204
n=8 graphs=11117 failures=0 min(M+1-F)=-0.000000 min(R-mu)=0.895751
n=9 part 0/4 graphs=65533 failures=0 min(M+1-F)=-0.000000 min(R-mu)=1.050649
n=9 part 1/4 graphs=47296 failures=0 min(M+1-F)=-0.000000 min(R-mu)=1.103389
n=9 part 2/4 graphs=94487 failures=0 min(M+1-F)=0.000000 min(R-mu)=1.080880
n=9 part 3/4 graphs=53764 failures=0 min(M+1-F)=0.333333 min(R-mu)=1.498261
\end{lstlisting}

\subsection{Direct check for order ten (Remark~\ref{rem:checks}(1))}\label{app:direct}
\begin{lstlisting}[language=C]
/* Independent direct check of Theorem 1.1 on all connected graphs of order n<=10 (graph6 on stdin),
   written separately from wc_small.py:
   F(G) = max over supports S (K_S nonsingular) of 1/sum(K_S^{-1} 1) when K_S^{-1} 1 > 0   (and 2 for |S|=1).
   M(G) exact as integer over DEN.  Reports min of M+1-F and the number of graphs with slack < 1e-9;
   the graphs with slack < 1e-6 are written to stderr for an exact re-check (exact_tight.py).           */
#include <stdio.h>
#include <string.h>
#include <math.h>
#define DEN 27720L
int main(void){
  char line[256]; long cnt=0, tight=0, fail=0; double minslack=1e9; int n=0;
  while(fgets(line,sizeof line,stdin)){
    line[strcspn(line,"\n")]=0; if(!line[0]) continue; cnt++;
    n=line[0]-63; int A[10][10]; memset(A,0,sizeof A);
    { int k=0; for(int j=1;j<n;j++) for(int i=0;i<j;i++){ int byte=line[1+k/6]-63; if((byte>>(5-k%6))&1) A[i][j]=A[j][i]=1; k++; } }
    int D[10][10];
    for(int i=0;i<n;i++) for(int j=0;j<n;j++) D[i][j]= i==j?0:(A[i][j]?1:99);
    for(int k=0;k<n;k++) for(int i=0;i<n;i++) for(int j=0;j<n;j++) if(D[i][k]+D[k][j]<D[i][j]) D[i][j]=D[i][k]+D[k][j];
    int deg[10]; long M=0;
    for(int i=0;i<n;i++){deg[i]=0; for(int j=0;j<n;j++) deg[i]+=A[i][j];}
    for(int i=0;i<n;i++) for(int j=i+1;j<n;j++) if(A[i][j]) M+=DEN/(deg[i]>deg[j]?deg[i]:deg[j]);
    double best=2.0;
    for(unsigned S=1;S<(1u<<n);S++){
      int idx[10], r=0; for(int i=0;i<n;i++) if(S>>i&1) idx[r++]=i;
      if(r<2) continue;
      double a[10][11];
      for(int i=0;i<r;i++){ for(int j=0;j<r;j++) a[i][j]= i==j?2.0:(double)D[idx[i]][idx[j]]; a[i][r]=1.0; }
      int sing=0;
      for(int c=0;c<r&&!sing;c++){
        int piv=c; double mx=fabs(a[c][c]);
        for(int i=c+1;i<r;i++) if(fabs(a[i][c])>mx){mx=fabs(a[i][c]); piv=i;}
        if(mx<1e-9){sing=1; break;}
        if(piv!=c) for(int j=c;j<=r;j++){double t=a[c][j]; a[c][j]=a[piv][j]; a[piv][j]=t;}
        for(int i=c+1;i<r;i++){ double f=a[i][c]/a[c][c]; if(f!=0.0) for(int j=c;j<=r;j++) a[i][j]-=f*a[c][j]; }
      }
      if(sing) continue;
      double w[10]; int pos=1; double tot=0;
      for(int i=r-1;i>=0;i--){ double s=a[i][r]; for(int j=i+1;j<r;j++) s-=a[i][j]*w[j]; w[i]=s/a[i][i]; }
      for(int i=0;i<r;i++){ if(w[i]<=1e-12) pos=0; tot+=w[i]; }
      if(!pos||tot<=0) continue;
      double val=1.0/tot; if(val>best) best=val;
    }
    double slack=(double)M/DEN+1.0-best;
    if(slack<minslack) minslack=slack;
    if(slack<1e-9) tight++;
    if(slack<1e-6) fprintf(stderr,"%s\n",line);
    if(slack<-1e-9){ fail++; printf("FAIL %s slack=%.9f\n",line,slack); }
  }
  printf("n=%d graphs=%ld failures=%ld min(M+1-F)=%.9f tight=%ld\n",n,cnt,fail,minslack,tight);
  return 0;
}
\end{lstlisting}
Output for the sixteen parts \texttt{nauty-geng -c -q 10 k/16}, followed by the output of the exact re-check (\texttt{extra/exact\_tight.py}) of the graphs written to the error stream:
\begin{lstlisting}
part 0/16: n=10 graphs=602072 failures=0 min(M+1-F)=0.000000000 tight=51
part 1/16: n=10 graphs=508270 failures=0 min(M+1-F)=0.000000000 tight=26
part 2/16: n=10 graphs=562497 failures=0 min(M+1-F)=0.000000000 tight=4
part 3/16: n=10 graphs=607709 failures=0 min(M+1-F)=0.500000000 tight=0
part 4/16: n=10 graphs=319966 failures=0 min(M+1-F)=0.833333333 tight=0
part 5/16: n=10 graphs=586883 failures=0 min(M+1-F)=0.333333333 tight=0
part 6/16: n=10 graphs=879635 failures=0 min(M+1-F)=0.816666667 tight=0
part 7/16: n=10 graphs=597745 failures=0 min(M+1-F)=0.000000000 tight=4
part 8/16: n=10 graphs=674677 failures=0 min(M+1-F)=0.166666667 tight=0
part 9/16: n=10 graphs=823542 failures=0 min(M+1-F)=0.000000000 tight=1
part 10/16: n=10 graphs=810514 failures=0 min(M+1-F)=0.464285714 tight=0
part 11/16: n=10 graphs=1078725 failures=0 min(M+1-F)=0.983333333 tight=0
part 12/16: n=10 graphs=1335009 failures=0 min(M+1-F)=0.583333333 tight=0
part 13/16: n=10 graphs=789448 failures=0 min(M+1-F)=0.666666667 tight=0
part 14/16: n=10 graphs=764578 failures=0 min(M+1-F)=0.500000000 tight=0
part 15/16: n=10 graphs=775301 failures=0 min(M+1-F)=1.233333333 tight=0
order 10: graphs re-checked exactly: 86; F = M+1 exactly: 86; F > M+1: 0
\end{lstlisting}


\subsection{Large irreducible decorated blocks (Remark~\ref{rem:checks}(5))}\label{app:big}
\begin{lstlisting}[language=Python]
# Stress test of Lemmas 4.1-4.5, inequality (6.1), the bound rho <= 1.3+1.43y and Theta > 1 on LARGE
# irreducible decorated blocks (up to about 100 sites); everything is computed from the definitions
# with exact rationals (irreducibility: M(G-e) > M(G) for every core edge e).
import sys, random, itertools, math
from fractions import Fraction as Fr
import networkx as nx

def Mval(G):
    d = dict(G.degree())
    return sum(Fr(1, max(d[u], d[v])) for u, v in G.edges())

def build(H, loaded):
    G = H.copy()
    for v in loaded: G.add_edge(v, ('leaf', v))
    return G

def irreducible(H, G):
    M = Mval(G)
    for u, v in H.edges():
        G.remove_edge(u, v)
        ok = Mval(G) > M
        G.add_edge(u, v)
        if not ok: return False
    return True

def analyse(H, loaded):
    """returns None if not an irreducible decorated block, else dict of checks"""
    loaded = set(loaded)
    if H.number_of_nodes() < 3 or not nx.is_biconnected(H): return None
    G = build(H, loaded)
    if not irreducible(H, G): return None
    d = dict(G.degree()); M = Mval(G)
    Mcore = sum(Fr(1, max(d[u], d[v])) for u, v in H.edges())
    y = math.sqrt(float(M) - 1)
    # Lemma 4.1 (P1)-(P3), from the definitions
    g = {v: d[v]*(d[v]-1) for v in H}
    ell = {v: sum(1 for w in G[v] if d[w] < d[v]) for v in H}
    a = {v: Fr(ell[v], g[v]) for v in H}; b_ = {v: Fr(d[v]-ell[v], g[v]) for v in H}
    for v in H:
        assert ell[v] >= 1 and d[v] >= 3
        if d[v] == 3: assert v in loaded and all(d[w] >= 4 for w in H[v])
    for u, v in H.edges():
        if d[u] > d[v]: u, v = v, u
        if d[u] < d[v]: assert a[u] > b_[v]
        else: assert a[u] >= b_[v] and a[v] >= b_[u]
    X = {}; xi = {}; gam = {}; theta = 0.0; worst_rho = -1e9; i2 = True
    for v in H.nodes():
        G2 = G.copy(); G2.remove_node(v)
        if v in loaded: G2.remove_node(('leaf', v))
        X[v] = M - Mval(G2)
        # Lemma 4.2: closed formula for X(v)
        form = Fr(1, d[v]) if v in loaded else Fr(0)
        for w in H[v]:
            form += b_[w] if d[w] > d[v] else (b_[w] - Fr(1, g[w]) if d[w] == d[v] else Fr(1, d[v]) - a[w])
        nb = list(H[v])
        form -= sum(Fr(1, g[w]) for i, w in enumerate(nb) for z in nb[i+1:] if H.has_edge(w, z) and d[w] == d[z])
        assert form == X[v]
        x = max(X[v], Fr(0))
        nl = sum(1 for w in H[v] if d[w] < d[v]); ne = sum(1 for w in H[v] if d[w] == d[v]); nh = sum(1 for w in H[v] if d[w] > d[v])
        peak = (ne == 0 and nh == 0)
        lam = Fr(0)
        if v in loaded:
            G3 = G.copy(); G3.remove_node(('leaf', v)); b = M - Mval(G3)
            assert b == Fr(d[v] - (nl + 1), d[v]*(d[v]-1))
            x = max(x, b)
            U = Fr((nl + 1)*(d[v]-1) + nh, d[v]*(d[v]-1))
            if X[v] > U: i2 = False
            if not (peak and d[v] >= 5): lam = Fr(1, d[v])
            # range claims of Lemma 4.4
            if d[v] == 3: assert Fr(1,3) <= x <= Fr(2,3)
            elif not peak: assert Fr(1, d[v]*(d[v]-1)) <= x <= 1 - Fr(d[v]-2, d[v]*(d[v]-1))
            elif d[v] == 4: assert x == Fr(1,2)
        assert X[v] <= 1
        xi[v] = x; gam[v] = Fr(5,4)*x + lam
        if x > 0:
            s = math.sqrt(float(x)); th = s/(y+s); theta += th
            worst_rho = max(worst_rho, float(gam[v])/th - (1.3 + 1.43*y))
        else:
            assert lam == 0
    h = H.number_of_nodes()
    f = min(Fr(t,2) + Fr(h-t, t-1) for t in range(2, h+1))
    return dict(h=h, M=float(M), I1=sum(X.values()) >= Mcore, I2=i2, I3=M <= sum(gam.values()),
                I4=worst_rho <= 1e-9, Theta=theta, count=M >= f, worst=worst_rho)

def hubs_and_middles(rng):
    k = rng.randint(2, 7)                       # hubs
    Q = nx.gnp_random_graph(k, rng.random(), seed=rng.randint(0, 10**9))
    H = nx.Graph(); H.add_nodes_from(range(k)); H.add_edges_from(Q.edges())
    nxt = k; loaded = set()
    pairs = list(itertools.combinations(range(k), 2))
    for (a, b) in pairs:
        m = rng.choice([0, 0, 1, 2, 3, 4, 6, 8]) if k > 2 else rng.randint(3, 30)
        for _ in range(m):
            H.add_edge(nxt, a); H.add_edge(nxt, b); loaded.add(nxt); nxt += 1
    # a few middles of core degree 3
    for _ in range(rng.randint(0, 4)):
        if k >= 3:
            tri = rng.sample(range(k), 3)
            for a in tri: H.add_edge(nxt, a)
            if rng.random() < 0.7: loaded.add(nxt)
            nxt += 1
    for a in range(k):
        if rng.random() < 0.4: loaded.add(a)
    return H, loaded

def multipartite(rng):
    k = rng.randint(2, 4); parts = [rng.randint(1, 9) for _ in range(k)]
    H = nx.complete_multipartite_graph(*parts)
    loaded = set()
    for v, dat in H.nodes(data=True):
        pass
    part_of = nx.get_node_attributes(H, 'subset')
    flags = [rng.random() < 0.5 for _ in range(k)]
    loaded = {v for v in H.nodes() if flags[part_of[v]]}
    return H, loaded

if __name__ == '__main__':
    seed = int(sys.argv[1]); N = int(sys.argv[2]); rng = random.Random(seed)
    found = 0; bad = 0; minT = 1e9; maxh = 0; hs = {}
    for it in range(N):
        H, loaded = (hubs_and_middles if rng.random() < 0.8 else multipartite)(rng)
        if H.number_of_nodes() < 4: continue
        r = analyse(H, loaded)
        if r is None: continue
        found += 1; maxh = max(maxh, r['h']); hs[r['h']] = hs.get(r['h'], 0) + 1
        minT = min(minT, r['Theta'])
        if not (r['I1'] and r['I2'] and r['I3'] and r['I4'] and r['count'] and r['Theta'] > 1):
            bad += 1; print('VIOLATION', r, sorted(H.edges()), sorted(loaded))
    print(f'seed={seed} trials={N} irreducible_found={found} max_h={maxh} violations={bad} minTheta={minT:.4f}')
    print('h histogram:', sorted(hs.items()))
\end{lstlisting}
Output:
\begin{lstlisting}
seed=1 trials=3000 irreducible_found=1376 max_h=104 violations=0 minTheta=1.3846
seed=2 trials=3000 irreducible_found=1405 max_h=111 violations=0 minTheta=1.3846
seed=3 trials=3000 irreducible_found=1476 max_h=106 violations=0 minTheta=1.3846
seed=4 trials=3000 irreducible_found=1383 max_h=108 violations=0 minTheta=1.3846
\end{lstlisting}

\end{document}
